\documentclass[11pt,letterpaper]{article}
\usepackage[T1]{fontenc}
\usepackage[utf8]{inputenc}
\usepackage{lmodern,amsmath,amssymb,amsthm,mathtools,bm}
\usepackage[margin=1in]{geometry}
\usepackage{microtype,booktabs,longtable,array,enumitem,xcolor}
\usepackage{titlesec,fancyhdr}
\usepackage[colorlinks=true,linkcolor=blue!45!black,citecolor=teal!60!black,urlcolor=blue!45!black,breaklinks=true]{hyperref}
\usepackage[nameinlink,noabbrev]{cleveref}
\definecolor{ink}{RGB}{25,47,67}
\definecolor{accent}{RGB}{0,105,111}
\titleformat{\section}{\large\bfseries\color{ink}}{\thesection}{.7em}{}
\titleformat{\subsection}{\normalsize\bfseries\color{accent}}{\thesubsection}{.7em}{}
\newtheorem{theorem}{Theorem}[section]
\newtheorem{lemma}[theorem]{Lemma}
\newtheorem{proposition}[theorem]{Proposition}
\newtheorem{corollary}[theorem]{Corollary}
\theoremstyle{definition}
\newtheorem{definition}[theorem]{Definition}
\newtheorem{example}[theorem]{Example}
\theoremstyle{remark}
\newtheorem{remark}[theorem]{Remark}
\DeclareMathOperator{\Li}{Li}
\DeclareMathOperator{\FP}{FP}
\DeclareMathOperator{\CT}{CT}
\DeclareMathOperator{\Res}{Res}
\DeclareMathOperator{\Rea}{Re}
\DeclareMathOperator{\Ima}{Im}
\newcommand{\Q}{\mathbb Q}
\newcommand{\Z}{\mathbb Z}
\newcommand{\C}{\mathbb C}
\newcommand{\N}{\mathbb N}
\newcommand{\dd}{\,\mathrm d}
\newcommand{\fall}[2]{#1^{\underline{#2}}}
\newcommand{\rising}[2]{(#1)_{#2}}
\newcommand{\repoSHA}{\texttt{446dca9bd65c20e409d1d2561d4ffad408103e2a}}
\setlength{\headheight}{14pt}
\pagestyle{fancy}
\fancyhf{}
\fancyhead[L]{\small\color{ink}Polygamma reconstruction and resonance}
\fancyhead[R]{\small ProveIt research continuation}
\fancyfoot[C]{\thepage}
\renewcommand{\headrulewidth}{.3pt}
\setlist{itemsep=.3em,topsep=.4em}
\setlength{\emergencystretch}{2em}
\numberwithin{equation}{section}
\allowdisplaybreaks[2]
\hypersetup{pdftitle={Exact Reconstruction and Higher Resonances in Polygamma--Zeta Calculus},pdfauthor={ProveIt research continuation}}
\title{\color{ink}\textbf{Exact Reconstruction and Higher Resonances\\in Polygamma--Zeta Calculus}\\[.6em]
\large Explicit even-zeta extraction, cubic tail moments,\\Gamma determinants, and spectral logarithmic integrals}
\author{Research continuation prepared for the ProveIt project}
\date{11 October 2026 (UTC)}
\begin{document}
\maketitle
\begin{abstract}
This article continues the polylogarithm, harmonic-tail, Gamma, and
Stieltjes programme at a fixed repository revision. Its main constructive
result recovers every even-zeta coordinate from the optimal moment prefix
for the pair of tail orders $(2,2b)$, with no matrix inverse or unspecified
minor left in the formula. The proof constructs an explicit Bernoulli
reproducing kernel and evaluates its rank-one perturbation. The resulting
polynomial weights give ordinary absolutely convergent polygamma series
and finite polylogarithmic boundary identities. Further developments concern
the optimal-prefix conjecture at tail order eight, all regular centered
cubic trigamma moments, the next two-particle Gamma-determinant
resonance, and spatial Mellin formulas for every spectral coefficient.
At order eight an exact symbolic determinant certificate and a uniform
Bernoulli congruence prove the conjecture outside one residue class
modulo sixteen. The cubic moments close in four fixed constant
coordinates. The resonance yields explicit generalized Stieltjes
identities, shift derivatives, and an anchored Hurwitz-zeta primitive.
All singular subtractions and coordinate conventions are stated explicitly.
Exact finite calculations and independent numerical diagnostics accompany
the proofs. The manuscript's short Gaussian $S_6$ and revised $S_8$ period
identities remain unresolved; no numerical independence claim is made.
\end{abstract}
\tableofcontents
\newpage
\section{Research target, source boundary, and conventions}
\label{sec:scope}

\subsection{What is continued here}

The source for this investigation is the ProveIt manuscript and its six
current incoming reports at commit
\begin{center}\small\repoSHA.\end{center}
The source was inspected on 11 October 2026. The canonical manuscript
already develops polylogarithmic identities, harmonic generating
functions, Gamma correlations, and several finite-part calculi
\cite{repo}. The incoming Bernoulli--Tornheim--digamma report proves
optimal-prefix rank for the unequal even tail-order pairs $(2,q)$,
$(4,q)$, and $(6,q)$, and asks for every reconstruction row in closed
form \cite[research question R2]{incoming-bernoulli}. The first rational
row is already known. This paper supplies the full inverse for $(2,q)$
and turns it into a family of evaluated ordinary series.

The six incoming archives were also inspected to distinguish new work
from statements already supplied there. In particular, the general
ordered-partition expansion of a product of tails, the symmetric cubic
descent, the finite differential-polynomial integration-by-parts
algorithm, and the first Gauss-determinant resonance are prior parts of
this programme \cite{incoming-cubic,incoming-differential,
incoming-determinantal,incoming-tail,incoming-multishift}.
They are not relabelled as new discoveries in this article.

The word ``new'' below refers to the advancement beyond these inspected
sources. A search of the relevant primary literature was used to identify
classical ingredients; it does not establish global historical priority.
The Bernoulli Hankel and orthogonal-polynomial background is classical
\cite{dilcher-jiu,dlmf-op}. We provide the particular algebraic proofs
needed here because the exceptional zeroth Bernoulli moment and the
tail matrix's omitted binomial terms affect the answer.

\subsection{Three different meanings of an evaluation}

We distinguish an ordinary convergent sum, a meromorphically continued
value, and a constant term of a specified cutoff. Where two of them
coincide, that coincidence is proved. In the centered two-tail problem,
the identities $\zeta(-2j,\tfrac12)=0$ make the requisite polynomial
subtractions have zero continued sum. That special cancellation does
not hold after an arbitrary change of centering.

The numerical checks in the accompanying files are diagnostics, not
interval certificates. An identity in a polynomial ring over $\Q$ is
checked exactly. A theorem about an infinite family or a continuum of
parameters is justified by the written proof. No result in the paper
uses a failed integer-relation search as an independence argument.

\subsection{Notation}

The Gamma function is $\Gamma$, the digamma function is
$\psi=\Gamma'/\Gamma$, and $\psi^{(m)}$ is its $m$th derivative.
For $p>1$ and $n\geq0$, write
\begin{equation}\label{eq:tails}
 H_n^{(p)}=\sum_{j=1}^n j^{-p},\qquad
 A_p(n)=\zeta(p)-H_n^{(p)}=\zeta(p,n+1),\qquad
 x_n=n+\tfrac12.
\end{equation}
For integer $p\geq2$,
\begin{equation}\label{eq:tail-polygamma}
 A_p(n)=\frac{(-1)^p}{(p-1)!}\psi^{(p-1)}(n+1).
\end{equation}
The Hurwitz-zeta and polygamma conventions agree with
\cite{dlmf-hurwitz,dlmf-polygamma}.
For $a>0$, the generalized Stieltjes constants are normalized by
\begin{equation}\label{eq:stieltjes-convention}
 \zeta(1+u,a)=\frac1u+
   \sum_{j=0}^{\infty}\frac{(-1)^j}{j!}\gamma_j(a)u^j.
\end{equation}
Thus $\gamma_0(a)=-\psi(a)$ and $\gamma_j=\gamma_j(1)$.

The Bernoulli convention is
\begin{equation}\label{eq:bernoulli-convention}
 \frac{te^{xt}}{e^t-1}=\sum_{j=0}^{\infty}B_j(x)\frac{t^j}{j!},
 \qquad B_j=B_j(0),\qquad B_1=-\tfrac12.
\end{equation}
The notation $(a)_j$ means a rising factorial. A falling factorial is
$a^{\underline j}=a(a-1)\cdots(a-j+1)$. Empty products are one.

For $|z|<1$, our depth-two polylogarithm is
\begin{equation}\label{eq:depth-two-convention}
 \Li_{p,q}(z)=\sum_{n>m\geq1}\frac{z^n}{n^p m^q}.
\end{equation}
The color belongs to the larger index. This convention matters in the
generating-function identities. At convergent endpoints we write
$\zeta_{\mathrm{MZ}}(p,q)=\sum_{n>m\geq1}n^{-p}m^{-q}$ for a
multiple zeta value; this differs from the Hurwitz function $\zeta(s,a)$.
The spectral double series $Z(s,t)$ in Section~\ref{mellin:section}
uses the same strict ordering on its stated convergence domain.
Auxiliary coefficients such as $C_j$ and $R_j$ are defined locally in
each family and do not identify constants from different sections.
All functions on the open unit disc are
defined by their series; real boundary limits are taken with $z\uparrow1$.
Complex powers of positive real numbers use the real logarithm. Additional
complex-parameter domains are stated locally where they are used.

\section{The exact tail matrix and its regularization}
\label{sec:tail-foundation}

Fix an integer $b\geq2$. Consider the Dirichlet series
\begin{equation}\label{eq:Db}
 D_b(s)=\sum_{n=0}^{\infty}A_2(n)A_{2b}(n)x_n^{-s},
 \qquad \Rea s>1-2b.
\end{equation}
Its centered moments are denoted by $D_{b,r}=D_b(-2r)$, $r\geq0$.
At a point outside the displayed convergence half-plane, this notation
means meromorphic continuation, not the original divergent series.

\subsection{Why every even moment is regular}

Define rational numbers
\begin{equation}\label{eq:tail-coefficients}
 t_{p,0}=\frac1{p-1},\qquad
 t_{p,j}=\frac{B_{2j}(1/2)(p)_{2j-1}}{(2j)!}\quad(j\geq1),
 \qquad
 a_{b,j}=\sum_{i=0}^j t_{2,i}t_{2b,j-i}.
\end{equation}
Euler--Maclaurin summation in the midpoint coordinate gives, to any fixed
order and with its usual power remainder,
\begin{align}
 A_p(n)&\sim\sum_{j\geq0}t_{p,j}x_n^{1-p-2j},
 \label{eq:tail-expansion}\\
 A_2(n)A_{2b}(n)&\sim\sum_{j\geq0}a_{b,j}x_n^{-2b-2j}.
 \label{eq:tail-product-expansion}
\end{align}
Only even corrections occur, since $B_{2j+1}(1/2)=0$.
The first two product coefficients are
\begin{equation}\label{eq:first-tail-coeff}
 a_{b,0}=\frac1{2b-1},\qquad
 a_{b,1}=-\frac{2b^2-b+1}{12(2b-1)}.
\end{equation}

\begin{proposition}[Ordinary subtraction at every even moment]
\label{prop:ordinary-subtraction}
The only possible poles of $D_b(s)$ are at
$s=1-2b-2j$, $j\geq0$. In particular every $D_{b,r}$ is regular.
For an even polynomial $W(x)=\sum_r w_rx^{2r}$, put
\begin{equation}\label{eq:subtraction-polynomial}
 R_{b,W}(x)=
 \sum_r w_r\sum_{j=0}^{r-b}a_{b,j}x^{2r-2b-2j},
\end{equation}
where an inner sum with negative upper limit is zero. Then
\begin{equation}\label{eq:ordinary-subtraction}
 \sum_{n=0}^{\infty}
  \{W(x_n)A_2(n)A_{2b}(n)-R_{b,W}(x_n)\}
 =\sum_r w_rD_{b,r}.
\end{equation}
The left side is an ordinary absolutely convergent series.
\end{proposition}
\begin{proof}
Subtract the first $J$ terms in \eqref{eq:tail-product-expansion} from
\eqref{eq:Db} and add back
$\sum_{j<J}a_{b,j}\zeta(s+2b+2j,\tfrac12)$.
The subtracted series converges normally in successively larger left
half-planes as $J$ increases. This gives the continuation and the list
of possible poles. None is an even integer.

Apply the same subtraction to each monomial of $W$. After removing
\eqref{eq:subtraction-polynomial}, the summand is $O(x_n^{-2})$.
The added-back values are all $\zeta(-2\ell,\tfrac12)$, $\ell\geq0$,
and each is zero because
$\zeta(s,\tfrac12)=(2^s-1)\zeta(s)$.
The evaluation of the meromorphic continuation is therefore exactly
the displayed ordinary sum. The finite number of monomials permits
termwise addition without an interchange issue.
\end{proof}

\subsection{The inherited Bernoulli coordinate transform}

Define transformed coordinates by
\begin{equation}\label{eq:C-from-D}
 C_j=\sum_{r=0}^j\binom{2j+1}{2r}2^{-2(j-r)}D_{b,r}.
\end{equation}
The inverse transform is
\begin{equation}\label{eq:D-from-C}
 D_{b,r}=\sum_{j=0}^r
   \frac{\binom{2r}{2j}}{2j+1}B_{2r-2j}(1/2)C_j.
\end{equation}
Both identities follow from the generating function
\eqref{eq:bernoulli-convention}. Indeed the exact prefix polynomial is
\begin{equation}\label{eq:prefix-polynomial}
 \sum_{n=0}^{N-1}(n+\tfrac12)^{2r}
 =\sum_{j=0}^r\frac{\binom{2r}{2j}}{2j+1}
       B_{2r-2j}(1/2)N^{2j+1}.
\end{equation}

The late coordinates $C_j$, $j\geq b$, close on ordinary even zeta
values. The following is the precise inherited reduction, in the
normalization used by the new inverse theorem
\cite[Sections 2--3]{incoming-bernoulli}.

\begin{proposition}[Late-coordinate formula]
\label{prop:late-coordinates}
Put $\mu_j=B_{2j}/(2j)!$ and $Z_k=\zeta(2k)$, including
$Z_0=-\tfrac12$. For $j\geq b$,
\begin{equation}\label{eq:late-coordinates}
 C_j=\sum_{k=0}^{b-1}
 \left\{\frac{(2j-1)!}{(2b-2k)!}\mu_{j-b+k}
 +\mathbf1_{k=0}\frac{2j-2b+1}{2}B_{2j-2b}\right\}Z_k.
\end{equation}
Consequently, if $\mathbf C=(C_b,\ldots,C_{2b-1})^{\mathsf T}$ and
$\mathbf Z=(Z_0,\ldots,Z_{b-1})^{\mathsf T}$, then
\begin{equation}\label{eq:M-system}
 \mathbf C=M_b\mathbf Z,
 \qquad
 (M_b)_{tk}=\frac{(2b+2t-1)!}{(2b-2k)!}\mu_{t+k}
       +\mathbf1_{k=0}\frac{2t+1}{2}B_{2t}.
\end{equation}
\end{proposition}

We recall the finite argument to make the boundary convention explicit.
For every odd integer $d>0$, summation by parts uses the exact prefix
expression
\begin{align}\label{eq:E-d}
 E_d(N)={}&N^dA_2(N)A_{2b}(N)
 +\sum_{k=1}^N k^{d-2b}A_2(k)
 +\sum_{k=1}^N k^{d-2}A_{2b}(k)
 +\sum_{k=1}^N k^{d-2b-2}.
\end{align}
It equals the moment sum whose prefix polynomial is $N^d$.
For a nonnegative power $m$, the finite identity
\begin{equation}\label{eq:faulhaber-tail}
 \sum_{k=1}^N k^m A_p(k)
 =A_p(N)F_m(N)+\sum_{\ell=1}^N\frac{F_m(\ell-1)}{\ell^p},
 \quad F_m(N)=\frac{B_{m+1}(N+1)-B_{m+1}(1)}{m+1},
\end{equation}
is obtained by reversing the finite sums, with the infinite part
$A_p(N)$ left intact. Substituting it into \eqref{eq:E-d}, expanding
the finite Bernoulli polynomials, and extracting the constant term
gives, for $d=2j+1\geq2b+1$, the two coefficients
\[
 B_{2j-2b+2k}
 \left\{
 \frac{\binom{2j-2b+2}{2-2k}}{2j-2b+2}
 +\frac{\binom{2j}{2b-2k}}{2j}\right\}.
\]
Here a binomial summand is omitted unless its lower index is at
least two and at most its upper index. The first term consequently
occurs only at $k=0$. Cancelling the factorials gives
\eqref{eq:late-coordinates}. At the transition $j=b$, the two
$-\gamma/2$ terms arising from the tail sums cancel the $+\gamma$
from the final harmonic sum in \eqref{eq:E-d}. This is why there is
no extra Euler constant in the zeroth column. This finite computation
is also the source proof of the proposition; it precedes the new
reconstruction problem.

The binomial omission rule is essential. Adding a lower-index-zero
term would change $M_b$. The matrix in \eqref{eq:M-system}, including
its exceptional zeroth column, is the one inverted below.

\subsection{Formal coordinates and the optimal prefix}

The source regards
\begin{equation}\label{eq:formal-list}
 C_0,\ldots,C_{b-1},\quad Z_0,\ldots,Z_{b-1}
\end{equation}
as $2b$ formal labels. The Bernoulli transform from $D_{b,0},\ldots,
D_{b,2b-1}$ is triangular, with the late block $M_b$. Thus
nonsingularity of $M_b$ gives full rank of this prefix.
It does not imply that the evaluated numbers in
\eqref{eq:formal-list} are arithmetically independent. In fact ordinary
even zeta values are rational multiples of powers of $\pi$.
Our inverse theorem solves a completely specified rational linear
problem on these formal labels and then evaluates the resulting
identities analytically.

\section{Every reconstruction row in explicit form}
\label{sec:inverse}

The main theorem of this section answers research question R2 of
\cite{incoming-bernoulli} for the complete $(2,2b)$ family. The answer
contains only factorials, Bernoulli numbers, and finite sums. In
particular, it does not leave an inverse matrix or an unknown
determinant in the definition of a reconstruction coefficient.

\subsection{A finite Pad\'e identity}

For $N\geq0$, define
\begin{equation}\label{eq:PN}
 P_N(t)=\sum_{k=0}^N
  \frac{(2N-k)!\,N!}{(2N)!\,k!\,(N-k)!}t^k.
\end{equation}
The exact initial Taylor coefficients of
$e^tP_N(-t)-P_N(t)$ vanish through degree $2N$, and
\begin{equation}\label{eq:Pade-remainder}
 e^tP_N(-t)-P_N(t)
 =\frac{(-1)^N(N!)^2}{(2N)!(2N+1)!}t^{2N+1}
   +O(t^{2N+2}).
\end{equation}
For completeness, the vanishing is the binomial identity
\[
 \sum_{k=0}^N(-1)^k\binom Nk(1+x)^{2N-k}
 =x^N(1+x)^N
\]
applied coefficient by coefficient. The first remaining coefficient is
\[
 \frac1{(2N)!}\sum_{k=0}^N
  \frac{(-1)^k\binom Nk}{2N+1-k}
 =\frac1{(2N)!}\int_0^1x^N(x-1)^N\dd x,
\]
which proves \eqref{eq:Pade-remainder} by the beta integral.

Set
\begin{equation}\label{eq:Fmu}
 F(z)=\sum_{j\geq0}\mu_jz^j
 =\frac{\sqrt z}{2}\coth\frac{\sqrt z}{2}.
\end{equation}
The right side is an even germ in $\sqrt z$, hence a single-valued
holomorphic germ in $z$. With $N=2n+1$, put
\[
 q_n(t^2)=\frac{P_{2n+1}(t)-P_{2n+1}(-t)}t,
 \qquad
 r_n(t^2)=\frac{P_{2n+1}(t)+P_{2n+1}(-t)}2.
\]
Then \eqref{eq:Pade-remainder} implies
\begin{equation}\label{eq:Fq-r}
 F(z)q_n(z)-r_n(z)=h_nz^{2n+1}+O(z^{2n+2}),
 \qquad
 h_n=\frac{((2n+1)!)^2}{(4n+2)!(4n+3)!}.
\end{equation}

\subsection{Two explicit orthogonal families}

Let $\mathcal L$ be the rational linear functional on $\Q[x]$ with
$\mathcal L(x^j)=\mu_j$. Define
\begin{align}
 \pi_n(x)&=\sum_{t=0}^n p_{nt}x^t,
 &p_{nt}&=
 \frac{2(2n+2t+1)!(2n+1)!}
 {(4n+2)!(2n-2t+1)!(2t)!},
 \label{eq:pi-coeff}\\
 c_n&=\frac{(2n+1)(2n+2)!}{(4n+2)!}.
 \label{eq:cn}
\end{align}
Thus $\pi_n(x)=x^nq_n(1/x)$ is monic. Coefficient comparison in
\eqref{eq:Fq-r} gives
\begin{equation}\label{eq:pi-moments}
 \mathcal L(\pi_n)=c_n,\qquad
 \mathcal L(x^k\pi_n)=0\ (1\leq k\leq n),\qquad
 \mathcal L(x^{n+1}\pi_n)=h_n.
\end{equation}
These are the inherited shifted-functional polynomials. The new inverse
construction requires orthogonality for $\mathcal L$ itself.

\begin{lemma}[The unshifted Bernoulli family]
\label{lem:Qorthogonal}
Put
\begin{align}
 Q_0(x)&=1,&\nu_0&=1,\notag\\
 Q_n(x)&=\pi_n(x)-\frac{c_n}{c_{n-1}}\pi_{n-1}(x),
 &\nu_n&=-\frac{c_n}{c_{n-1}}h_{n-1}\quad(n\geq1).
 \label{eq:Qnu}
\end{align}
Then the $Q_n$ are monic, pairwise orthogonal for $\mathcal L$, and
$\mathcal L(Q_n^2)=\nu_n\ne0$. Write
$q_{nt}=[x^t]Q_n(x)$, with $q_{nt}=0$ if $t>n$.
\end{lemma}
\begin{proof}
The zeroth moment of $Q_n$ vanishes by \eqref{eq:pi-moments}.
For $1\leq k<n$, both terms in
$\mathcal L(x^kQ_n)$ vanish by the same formula. Therefore $Q_n$
is orthogonal to every polynomial of degree less than $n$.
As $Q_n$ is monic,
\[
 \mathcal L(Q_n^2)=\mathcal L(x^nQ_n)
 =-\frac{c_n}{c_{n-1}}\mathcal L(x^n\pi_{n-1})
 =-\frac{c_n}{c_{n-1}}h_{n-1}.
\]
All factorials in $c_n,h_n$ are positive, so the norm is nonzero.
The functional is not positive: $\nu_0>0$ and $\nu_n<0$ for
$n\geq1$. The proof is algebraic and uses no positivity assumption.
\end{proof}

\begin{theorem}[Explicit inverse Bernoulli kernel]
\label{thm:kernel}
For $b\geq1$, let $H_b=(\mu_{t+k})_{0\leq t,k<b}$ and set
\begin{equation}\label{eq:kernel}
 K_b(x,y)=\sum_{n=0}^{b-1}\frac{Q_n(x)Q_n(y)}{\nu_n},
 \qquad
 K_{kt}^{(b)}=[x^ky^t]K_b(x,y)
 =\sum_{n=0}^{b-1}\frac{q_{nk}q_{nt}}{\nu_n}.
\end{equation}
Then $H_b$ is nonsingular and
$H_b^{-1}=(K_{kt}^{(b)})_{0\leq k,t<b}$.
Moreover,
\begin{equation}\label{eq:Kfirst}
 K_{0t}^{(b)}=\frac{p_{b-1,t}}{c_{b-1}}.
\end{equation}
\end{theorem}
\begin{proof}
Every polynomial $f$ of degree less than $b$ has the orthogonal
expansion
\[
 f(x)=\sum_{n=0}^{b-1}
   \frac{\mathcal L(fQ_n)}{\nu_n}Q_n(x).
\]
It follows that $\mathcal L_y(K_b(x,y)f(y))=f(x)$.
Taking $f(y)=y^j$ and comparing coefficients of $x^k$ gives
$\sum_tK_{kt}^{(b)}\mu_{t+j}=\delta_{kj}$, which proves the matrix
identity. The polynomial $\pi_{b-1}/c_{b-1}$ has moments one at
degree zero and zero at degrees $1,\ldots,b-1$ by
\eqref{eq:pi-moments}; uniqueness of the first inverse row proves
\eqref{eq:Kfirst}.
\end{proof}

\subsection{The complete mixed inverse}

Define the following finite rational data, for $0\leq t,k<b$:
\begin{align}
 u_t^{(b)}&=\mu_t
   \frac{(2b)!(2t+1)!}{2(2b+2t-1)!},
 \label{eq:ut}\\
 v_k^{(b)}&=\sum_{t=0}^{b-1}K_{kt}^{(b)}u_t^{(b)},
 &\delta_b&=2-B_{2b-2}.
 \label{eq:vkdelta}
\end{align}
Here every $K_{kt}^{(b)}$ is already given by the explicit finite sum
\eqref{eq:kernel} and the factorial coefficients
\eqref{eq:pi-coeff}--\eqref{eq:Qnu}.

\begin{theorem}[All reconstruction rows]
\label{thm:all-rows}
For every $b\geq2$, $\delta_b\ne0$, and
\begin{equation}\label{eq:all-rows}
 \boxed{\displaystyle
 \mathcal R_{b;k,t}=
 \frac{(2b-2k)!}{(2b+2t-1)!}
 \left(
 K_{kt}^{(b)}-
 \frac{v_k^{(b)}}{\delta_b}\frac{p_{b-1,t}}{c_{b-1}}
 \right)
 =(M_b^{-1})_{kt}.}
\end{equation}
Consequently every even-zeta coordinate is recovered by
\begin{equation}\label{eq:all-coordinates}
 \boxed{\displaystyle
 \zeta(2k)=\sum_{t=0}^{b-1}\mathcal R_{b;k,t}C_{b+t},
 \qquad 0\leq k<b.}
\end{equation}
This gives all rows in a common closed finite-sum formula.
\end{theorem}
\begin{proof}
Let
\[
 R_b=\operatorname{diag}_{t}(2b+2t-1)!,\qquad
 E_b=\operatorname{diag}_{k}\frac1{(2b-2k)!}.
\]
The exact entries in \eqref{eq:M-system} give
\begin{equation}\label{eq:rank-one}
 M_b=R_b(H_b+u^{(b)}e_0^{\mathsf T})E_b.
\end{equation}
By \eqref{eq:Kfirst}, factorial cancellation gives
\begin{align}\label{eq:v0}
 v_0^{(b)}
 &=\frac1{2b-1}\sum_{t=0}^{b-1}
   \binom{2b-1}{2t}(2t+1)B_{2t}
 =1-B_{2b-2}.
\end{align}
To prove the last equality directly, put $N=2b-1\geq3$ and use
\[
 \sum_{j=0}^N\binom NjB_j=B_N(1)=0,\qquad
 \sum_{j=0}^N j\binom NjB_j=-NB_{N-1}.
\]
After removing the $j=1$ terms and using the vanishing of the higher
odd Bernoulli numbers, their sum is exactly \eqref{eq:v0}.

We include the elementary nonvanishing argument. In the localization
$\Z_{(2)}$ of rationals with odd denominator, the Bernoulli recurrence
is
\[
 (2m+1)(2B_{2m})=2m-1-
  \sum_{j=1}^{m-1}\binom{2m+1}{2j}(2B_{2j}).
\]
Induction and
$\sum_{j=1}^{m-1}\binom{2m+1}{2j}=2^{2m}-2m-2$
show that $2B_{2m}$ is an odd element of $\Z_{(2)}$.
It cannot equal four. Hence $B_{2m}\ne2$, proving $\delta_b\ne0$.

Let $K=H_b^{-1}$ and $v=Ku$. A direct multiplication, using
$1+e_0^{\mathsf T}v=\delta_b$, shows
\[
 (H_b+ue_0^{\mathsf T})^{-1}
 =K-\frac{ve_0^{\mathsf T}K}{\delta_b}.
\]
Insert \eqref{eq:Kfirst} and multiply by the two inverse diagonal
factors from \eqref{eq:rank-one}. This is \eqref{eq:all-rows}.
Finally apply it to \eqref{eq:M-system}.
\end{proof}

\begin{remark}[Recovery of the earlier rational row]
The row $k=0$ simplifies to
\begin{equation}\label{eq:rational-row}
 \mathcal R_{b;0,t}=
 \frac{2}{(2b-1)(2-B_{2b-2})}\binom{2b-1}{2t}.
\end{equation}
Thus the previously known rational extraction is recovered exactly.
Rows $k=1,\ldots,b-1$ in \eqref{eq:all-rows} are the missing rows
asked for in the source. The first nontrivial complete inverse is
\begin{equation}\label{eq:b2inverse}
 M_2=
 \begin{pmatrix}3/4&1/4\\2/3&-1/12\end{pmatrix},\qquad
 M_2^{-1}=\frac1{11}
 \begin{pmatrix}4&12\\32&-36\end{pmatrix}.
\end{equation}
These small matrices are included as a transparent normalization check.
\end{remark}

\section{Evaluated ordinary series and polylogarithmic boundary identities}
\label{sec:filters}

\subsection{A polynomial for each prescribed even zeta value}

For $0\leq k<b$, define the odd polynomial and its centered difference
\begin{align}
 S_{b,k}(y)&=\sum_{t=0}^{b-1}\mathcal R_{b;k,t}y^{2b+2t+1},
 \label{eq:S-filter}\\
 W_{b,k}(x)&=S_{b,k}(x+\tfrac12)-S_{b,k}(x-\tfrac12),
 &R_{b,k}(x)&=R_{b,W_{b,k}}(x).
 \label{eq:W-filter}
\end{align}
Here $R_{b,W}$ is the explicit asymptotic polynomial in
\eqref{eq:subtraction-polynomial}. Thus $W_{b,k}$ is an even rational
polynomial of degree at most $4b-2$, and $R_{b,k}$ is even of degree
at most $2b-2$.

\begin{theorem}[All even-zeta extraction series]
\label{thm:extraction-series}
For every $b\geq2$ and every $0\leq k<b$,
\begin{equation}\label{eq:extraction-series}
 \boxed{\displaystyle
 \sum_{n=0}^{\infty}
  \left\{W_{b,k}(x_n)
    \bigl(\zeta(2)-H_n^{(2)}\bigr)
    \bigl(\zeta(2b)-H_n^{(2b)}\bigr)-R_{b,k}(x_n)\right\}
 =\zeta(2k).}
\end{equation}
The summand is $O(n^{-2})$, and the sum is ordinary and absolutely
convergent. Equivalently, its tail product may be replaced by
\begin{equation}\label{eq:extraction-polygamma}
 \frac{\psi'(n+1)\psi^{(2b-1)}(n+1)}{(2b-1)!}.
\end{equation}
Every coefficient in \eqref{eq:extraction-series} has the finite formula
\eqref{eq:all-rows}; the theorem is constructive at arbitrary order.
\end{theorem}
\begin{proof}
For an odd monomial,
\[
 (x+\tfrac12)^{2j+1}-(x-\tfrac12)^{2j+1}
 =\sum_{r=0}^j\binom{2j+1}{2r}2^{-2(j-r)}x^{2r}.
\]
Consequently the continued moment associated with $W_{b,k}$ is
$\sum_t\mathcal R_{b;k,t}C_{b+t}=\zeta(2k)$ by
\eqref{eq:C-from-D} and \cref{thm:all-rows}. Apply
\cref{prop:ordinary-subtraction} to turn that continued value into
the ordinary sum in \eqref{eq:extraction-series}. The polygamma
form follows from \eqref{eq:tail-polygamma}.
\end{proof}

\begin{proposition}[Uniqueness on the optimal formal prefix]
\label{prop:filter-unique}
Among even polynomials of degree at most $4b-2$, $W_{b,k}$ is the
unique one whose moment has formal coordinate vector $Z_k$ in
\eqref{eq:formal-list}.
\end{proposition}
\begin{proof}
The centered difference maps odd polynomials of degree at most
$4b-1$ bijectively onto even polynomials of degree at most $4b-2$:
in the monomial bases its matrix is triangular with nonzero diagonal.
Write the preimage of a candidate polynomial as
$\sum_{j=0}^{2b-1}s_jy^{2j+1}$. Its formal moment is
$\sum_js_jC_j$. The first $b$ labels in \eqref{eq:formal-list}
force $s_0=\cdots=s_{b-1}=0$. The remaining row must solve
$sM_b=e_k^{\mathsf T}$, whose unique solution is the row
\eqref{eq:all-rows}. This is a formal uniqueness statement; accidental
relations between evaluated constants are not excluded.
\end{proof}

\subsection{A compact trigamma--third-polygamma evaluation}

The row $k=1$ for $b=2$ gives a relatively short new member of the
family. With $x_n=n+\tfrac12$, put
\[
 P(x)=4032x^6+2480x^4-524x^2-23.
\]
Then
\begin{equation}\label{eq:explicit-zeta2-tail}
 \boxed{\displaystyle
 \sum_{n=0}^{\infty}
 \left\{64(63x_n^2+2)-3P(x_n)A_2(n)A_4(n)\right\}
 =528\zeta(2)=88\pi^2.}
\end{equation}
Both large terms in braces must be combined before summation. Each
separate sum diverges. Their combination is absolutely convergent.

Using $t_n=2n+1$ gives the equivalent integral-coefficient
polygamma identity
\begin{equation}\label{eq:explicit-zeta2-polygamma}
 \boxed{\begin{aligned}
 \sum_{n=0}^{\infty}\Bigl\{
 &32(63t_n^2+8)\\
 &-(63t_n^6+155t_n^4-131t_n^2-23)
          \psi'(n+1)\psi'''(n+1)\Bigr\}=176\pi^2.
 \end{aligned}}
\end{equation}
This is proved by the exact row identity, not by recognizing a decimal
value. Numerically, \eqref{eq:explicit-zeta2-tail} evaluates to
\[
 868.525187295863558457435207989101299907605548\ldots,
\]
as does $88\pi^2$.

For a direct coefficient check, the inverse row from
\eqref{eq:b2inverse} is $(32,-36)/11$ and
\[
 S_{2,1}(y)=\frac{32y^5-36y^7}{11}.
\]
Its centered difference is $-3P(x)/528$. The nonnegative-power
part of its product with the two tails is
$-64(63x^2+2)/528$, giving
\eqref{eq:explicit-zeta2-tail} immediately. This short calculation
also fixes the sign of the subtraction.

\subsection{An evaluation at the next tail order}

To illustrate recovery of a coordinate beyond $\zeta(2)$, let
\begin{align*}
 P_3(t)={}&2145t^{10}+40725t^8+162302t^6\\
          &+146110t^4+22161t^2+61,\\
 Q_3(t)={}&137280t^4+1874240t^2+1263872.
\end{align*}
The row $k=2$ for $b=3$ yields
\begin{equation}\label{eq:explicit-zeta4}
 \boxed{\displaystyle
 \sum_{n=0}^{\infty}
  \{Q_3(t_n)-5P_3(t_n)A_2(n)A_6(n)\}
 =19520\zeta(4).}
\end{equation}
Once again the combined summand is $O(n^{-2})$. The product is
$\psi'(n+1)\psi^{(5)}(n+1)/120$.
The supplied coefficient tables include all rows for $2\leq b\leq6$;
the generation code implements the theorem for any $b\geq2$.
Those tables are examples of an all-order formula, not the evidence
on which its all-order validity rests.

\subsection{Finite polylogarithmic generators}

Let $\theta=z\,d/dz$. For integers $p,q\geq2$, the ordinary
generating function of the tail product is
\begin{equation}\label{eq:Gpq}
 G_{p,q}(z)=\sum_{n=0}^{\infty}A_p(n)A_q(n)z^n
 =\frac{N_{p,q}(z)}{1-z},\qquad |z|<1,
\end{equation}
where
\begin{align}\label{eq:Npq}
 N_{p,q}(z)={}&\zeta(p)\zeta(q)-\zeta(q)\Li_p(z)
             -\zeta(p)\Li_q(z)\notag\\
 &+\Li_{p,q}(z)+\Li_{q,p}(z)+\Li_{p+q}(z).
\end{align}
This is the tail-product bridge already present in
\cite{incoming-tail}. A short direct proof is useful: multiply the
series by $1-z$, use
$A_p(n-1)=A_p(n)+n^{-p}$, and then use
\[
 \sum_{n\geq1}\frac{z^nH_n^{(q)}}{n^p}
   =\Li_{p,q}(z)+\Li_{p+q}(z).
\]
All series are absolutely convergent in the open disc, which justifies
each rearrangement.

\begin{theorem}[Evaluated polylogarithmic boundary operators]
\label{thm:polylog-filter}
For every $b\geq2$ and $0\leq k<b$, define
\begin{equation}\label{eq:Phi-filter}
 \Phi_{b,k}(z)=
 W_{b,k}(\theta+\tfrac12)\frac{N_{2,2b}(z)}{1-z}
 -R_{b,k}(\theta+\tfrac12)\frac1{1-z}.
\end{equation}
This is a finite differential expression in ordinary and depth-two
polylogarithms, holomorphic for $|z|<1$. It extends continuously to
the closed unit disc, and
\begin{equation}\label{eq:Phi-boundary}
 \boxed{\Phi_{b,k}(1)=\zeta(2k).}
\end{equation}
\end{theorem}
\begin{proof}
Termwise differentiation of normally convergent power series gives
\[
 \Phi_{b,k}(z)=\sum_{n=0}^{\infty}
 \{W_{b,k}(x_n)A_2(n)A_{2b}(n)-R_{b,k}(x_n)\}z^n.
\]
The coefficients belong to $\ell^1$ by
\cref{thm:extraction-series}. The Weierstrass test gives uniform
convergence on $|z|\leq1$, and the boundary value is exactly
\eqref{eq:extraction-series}. Individual differentiated terms in
\eqref{eq:Phi-filter} can diverge at $z=1$; the combined expression
has the proved continuous extension.
\end{proof}

For example, \eqref{eq:explicit-zeta2-tail} has the finite functional
form
\begin{equation}\label{eq:explicit-boundary-operator}
 \lim_{z\uparrow1}
 \left\{64\bigl(63(\theta+\tfrac12)^2+2\bigr)\frac1{1-z}
       -3P(\theta+\tfrac12)\frac{N_{2,4}(z)}{1-z}\right\}
 =88\pi^2.
\end{equation}
The theorem evaluates a prescribed cancellation of boundary jets.
It does not claim a reduction of arbitrary depth-two polylogarithms
to depth one. The exact operator and the correct boundary constant
are both part of the identity.

\section{Order eight: an infinite family beyond the proved small orders}
\label{rank8:section}

The incoming \emph{Bernoulli, Tornheim, and Digamma Identities} report
\cite[general small-determinant reduction and research question R1]{incoming-bernoulli}
proves the optimal-prefix conjecture when the smaller even order is
$2$, $4$, or $6$, and reduces every fixed smaller order to a finite
determinant.  The following result establishes a substantial part of
the first remaining order.  Its proof uses an exact symbolic
determinant certificate and a uniform Bernoulli congruence.

\begin{theorem}[The order-eight prefix outside one residue class]
\label{rank8:theorem}
Let $p=8$ and $q=2b>8$.  If
\begin{equation}\label{rank8:scope}
 b\not\equiv1\pmod {16},
\end{equation}
then the first $q$ centered moments have full formal rank $q$ on
the inherited coordinate list
$C_0,\ldots,C_{b-1},Z_0,\ldots,Z_{b-1}$.
The assertion concerns the coefficient matrix; it does not assert
arithmetic independence of the evaluated coordinates.

For $b\ge8$ satisfying \eqref{rank8:scope}, the normalized determinant
ratio $\rho_b$ defined below has the exact valuation
\begin{equation}\label{rank8:valuation}
 v_2(\rho_b)=v_2(b-1)-4.
\end{equation}
No all-$b$ valuation formula is asserted.
\end{theorem}

\subsection{The inherited reduction, with its normalization fixed}

Here the centered moments are the regular values
$D_{8,2b}(-2r)$, $0\le r<2b$, where
\[
 D_{8,2b}(s)=\sum_{n\ge0}A_8(n)A_{2b}(n)x_n^{-s}
 \quad(\Re s>-2b-5)
\]
is continued by the same midpoint subtraction as in
Proposition~\ref{prop:ordinary-subtraction}. The coordinates $C_j$
use \eqref{eq:C-from-D} with $D_{b,r}$ replaced by $D_{8,2b}(-2r)$.
The zeta coordinates are again $Z_k=\zeta(2k)$.

Set $\mu_r=B_{2r}/(2r)!$ for $r\ge0$ and $\mu_r=0$ for $r<0$,
and put
\[
 H_{4,b}=(\mu_{i+j-3})_{0\le i,j<b},\quad
 R_{4,b}=\operatorname{diag}_{0\le i<b}(2b+2i-7)!,\quad
 E_b=\operatorname{diag}_{0\le j<b}\frac1{(2b-2j)!}.
\]
The source's $b$-square late-coordinate matrix is explicitly
\begin{equation}\label{rank8:source-matrix}
 (M_{4,b})_{ik}=B_{2i+2k-6}
 \left\{\frac{\binom{2i+2}{8-2k}}{2i+2}
       +\frac{\binom{2b+2i-6}{2b-2k}}{2b+2i-6}\right\},
 \qquad 0\le i,k<b.
\end{equation}
An entry with negative Bernoulli index is zero. A binomial summand is
omitted unless its lower index is at least $2$ and at most its upper
index. Thus every matrix in the determinant ratio is specified. Define
\begin{equation}\label{rank8:rho}
 \rho_b=\frac{\det M_{4,b}}
 {\det R_{4,b}\det E_b\det H_{4,b}}.
\end{equation}
The inherited shifted-Hankel theorem makes every factor in the
denominator nonzero.  The first-$q$ determinant is
$\det M_{4,b}/(4b-1)!!$, so it suffices to prove $\rho_b\ne0$.

Here is a four-square formula whose size is independent of $b$.
For $0\le i,\ell\le3$ set
\begin{align*}
 n_i&=b-4+i,&N_{i\ell}&=2(n_i-\ell)+1,&R_\ell&=2b-6+2\ell,\\
 f_\ell(n)&=\frac1{(2\ell+1)!}
       \prod_{r=0}^{2\ell-1}(2n-r)(2n+3+r),&F_{i\ell}&=f_\ell(n_i).
\end{align*}
An empty product is $1$.  Define the polynomial weight
\[
 W_{i\ell}(z)=(z+2\ell+1)
                 (2n_i+2\ell+z+1)^{\underline{2i}},
\]
where the underline denotes a falling factorial, and its finite
Newton coefficients
\begin{equation}\label{rank8:newton}
 d_{i\ell h}=\frac1{h!}\sum_{j=0}^{h}
       (-1)^{h-j}\binom hj W_{i\ell}(N_{i\ell}-j).
\end{equation}
Introduce seven independent symbols $U_1,\ldots,U_7$.  Set
\begin{equation}\label{rank8:A}
 \begin{split}
 \mathcal A_{i\ell}={}&2F_{i\ell}+
 \frac{R_\ell!}
 {(2n_i+1)(2n_i+2)(2\ell+2)!N_{i\ell}!}\,
 \sum_{\substack{1\le h\le2i+1\\h\ {
m odd}}}
 d_{i\ell h}N_{i\ell}^{\underline h}
 U_{4-i+\ell+(h-1)/2}.
 \end{split}
\end{equation}
The ratio of factorials is interpreted as a rational function of $b$
by cancellation of finite products; its index difference
$R_\ell-N_{i\ell}=1+4\ell-2i$ is fixed.

For every integer $b\ge8$ the inherited small-determinant reduction
gives
\begin{equation}\label{rank8:ratio-small}
 \rho_b=
 \left.\frac{\det\mathcal A}{\det F}
 \right|_{U_j=2B_{2b-2j}}.
\end{equation}
To check the normalization, reverse the four source columns and
divide row $i$ by its nonzero Pad\'e constant $c_{n_i}$.  The
unperturbed entries become $F_{i\ell}$.  For the perturbation, apply
\eqref{rank8:newton} to the source weighted Bernoulli sum, using
the Newton identity
$W_{i\ell}(z)=\sum_h d_{i\ell h}(N_{i\ell}-z)^{\underline h}$
and
\[
 \sum_{j=0}^{N}\binom Nj B_j(N-j)^{\underline h}
      =N^{\underline h}B_{N-h}(1).
\]
When $b\ge8$, $N_{i\ell}-h\ge2$ for every retained odd $h$, and
the even $h$ terms have odd Bernoulli index at least $3$ and vanish.
The excluded $j=1$ term contributes exactly another $F_{i\ell}$.
This explains the coefficient $2F$ in \eqref{rank8:A} and the factor
$U_j=2B_{2b-2j}$; neither can be dropped.  At smaller $b$ the
exceptional $B_0$ and $B_1(1)$ terms require the direct finite
calculation recorded below.

\subsection{The exact polynomial certificate}

Put
\begin{align}
 D_8(b)={}&256b^8-3584b^7+18496b^6-40544b^5
 +25984b^4\notag\\
 &+18424b^3-5856b^2-8316b-2025.
 \label{rank8:D8}
\end{align}
This integer is odd at every integer $b$, since every nonconstant
coefficient is even.

\begin{lemma}[Finite integer certificate]\label{rank8:certificate}
The matrices \eqref{rank8:A} and $F$ satisfy
\begin{align}
 \det F={}&\frac{256}{4725}(b-2)(b-1)(2b-5)(2b-3)
 (4b-11)(4b-9)\notag\\
 &\hspace{15mm}\cdot(4b-7)^2(4b-5)(4b-3)D_8(b),
 \label{rank8:Fdet}\\
 \frac{\det\mathcal A}{\det F}
   ={}&\frac{P(b;U_1,\ldots,U_7)}{576(2b-5)D_8(b)},
 \label{rank8:certificate-ratio}
\end{align}
where $P\in\mathbb Z[b,U_1,\ldots,U_7]$.  It has $94$ monomials
in the seven $U$ variables, with integer polynomial coefficients in
$b$.  Under the indicated polynomial substitution, it obeys
\begin{equation}\label{rank8:congruence-certificate}
 \boxed{P\bigl(b;1+4(b-j)+16(b-j)^2\ (1\le j\le7)\bigr)
       \equiv16b^2+44b+4\pmod {64}.}
\end{equation}
The congruence is coefficientwise in $\mathbb Z[b]$.
\end{lemma}

\begin{proof}
This is a finite algebraic certificate. The accompanying
\texttt{rank\_eight\_certificate.json} gives
every coefficient of $P$.  The independent verifier rebuilds
\eqref{rank8:newton}--\eqref{rank8:A}, without reading a precomputed
matrix.  In each row it clears the least common denominator and
computes the determinant as the sum of its $24$ signed permutation
products in $\mathbb Z[b,U_1,\ldots,U_7]$.  Cross multiplication
verifies \eqref{rank8:Fdet} and \eqref{rank8:certificate-ratio} by
exact coefficient comparison.  It then performs the seven displayed
substitutions and checks that every coefficient of the difference in
\eqref{rank8:congruence-certificate} is divisible by $64$.

These instructions, together with the finite formulas defining the
entries, specify a reproducible exact certificate.  The calculation
uses neither interpolation from values of $b$ nor floating-point
determinants.  In particular, it certifies polynomial identities for
all parameters, rather than inferring them from a sample.
\end{proof}

\subsection{A uniform Bernoulli congruence}

Let $\mathbb Z_{(2)}$ denote the rational numbers with odd
denominator; congruences below are in this ring.

\begin{lemma}[Doubled Bernoulli numbers modulo $64$]
\label{rank8:bernoulli}
For every integer $m\ge4$,
\begin{equation}\label{rank8:bernoulli64}
 \boxed{2B_{2m}\equiv1+4m+16m^2\pmod {64}.}
\end{equation}
\end{lemma}

\begin{proof}
First, $2B_{2m}\in\mathbb Z_{(2)}$ is odd for every $m\ge1$.
This follows inductively from
\[
 (2m+1)(2B_{2m})
 =2m-1-\sum_{j=1}^{m-1}\binom{2m+1}{2j}(2B_{2j}).
\]
Indeed the sum of the binomial coefficients on the right is even,
and the denominator $2m+1$ is odd.  Thus all even Bernoulli
numbers have $2$-adic valuation at least $-1$, a fact also proved
in the incoming report.

Apply the Faulhaber formula to
$S_{2m}(8)=\sum_{a=0}^{7}a^{2m}$.  Its final term is $8B_{2m}$.
For $m\ge2$, every other nonzero term with an even Bernoulli index
contains at least $8^3$; its valuation is at least $8$ after the
Bernoulli denominator is included.  The $B_1$ and $B_0$ terms have
still greater valuation.  Division by $2m+1$ has no effect because
that integer is odd.  Consequently
\[
 2B_{2m}\equiv\frac14S_{2m}(8)\pmod {64}.
\]
When $m\ge4$, each even summand $a^{2m}/4$ vanishes modulo $64$.
For the four odd summands, expand the integer powers
$9^m=(1+8)^m$, $25^m=(1+24)^m$, and $49^m=(1+48)^m$.
After division by $4$, the first coefficients are
\[
 \frac{1+9^m+25^m+49^m}{4}
 \equiv1+20m+736\binom m2\pmod {64}.
\]
The coefficient of $\binom m3$ is $31232$, divisible by $64$;
all subsequent terms have valuation at least $10$.  Finally,
$736\equiv32\pmod {64}$, and
$20m+32\binom m2=4m+16m^2$.  This proves the lemma.
\end{proof}

\subsection{Proof of the order-eight theorem}

For $b\ge11$ all seven indices $m=b-j$ in
\eqref{rank8:ratio-small} are at least $4$.  Since $P$ has integer
coefficients, Lemma~\ref{rank8:bernoulli} may be substituted into
Lemma~\ref{rank8:certificate}, giving
\begin{equation}\label{rank8:P64}
 P(b;2B_{2b-2},\ldots,2B_{2b-14})
 \equiv16b^2+44b+4
 \equiv4(b-1)(4b+15)\pmod {64}.
\end{equation}
The last factor is odd.  If $b\not\equiv1\pmod {16}$, then
$v_2(b-1)\le3$, so the right side of \eqref{rank8:P64} has valuation
at most $5$, and
\[
 v_2(P)=v_2(b-1)+2.
\]
The denominator in \eqref{rank8:certificate-ratio} has valuation
$6$: $576=2^6\cdot9$, while $2b-5$ and $D_8(b)$ are odd.
This proves \eqref{rank8:valuation} and nonvanishing for all the
stated $b\ge11$.

The remaining six values $5\le b\le10$ are computed from the
original $b$-square source matrix, with every omitted binomial
summand treated exactly as prescribed.  Their nonzero ratios are
\begin{center}
\begin{tabular}{c r}
\hline
$b$ & $\rho_b$\\
\hline
5 & $309713/27640$\\
6 & $263589503/10455480$\\
7 & $-6698504467/244114290$\\
8 & $1826240201143/4860702000$\\
9 & $-334092867653747/85606130610$\\
10 & $711292958938006619/13975082830800$\\
\hline
\end{tabular}
\end{center}
At $b=8,9,10$ their valuations are $-4,-1,-4$, agreeing with
\eqref{rank8:valuation}.  The cases $b=5,6,7$ are asserted only
as nonzero, not as instances of that valuation formula.
This completes the proof of Theorem~\ref{rank8:theorem}.

\subsection{Limits of the result and an informative failed extrapolation}

The unevaluated class $b\equiv1\pmod {16}$ remains outside this
theorem.  Computed values in it do not prove its general nonvanishing.
Moreover, an apparent extension of \eqref{rank8:valuation}, suggested
by low-order experiments during this work, is false.  Independent
exact evaluation with ordinary rational Bernoulli numbers gives
\[
 \begin{array}{c|c|c|c}
 b&P\bmod 65536&v_2(\rho_b)&v_2(b-1)-4\\\hline
 65&52224&4&2\\
 129&24320&2&3
 \end{array}
\]
The residues have valuations $10$ and $8$, respectively; the
denominator still has valuation $6$.  These are counterexamples
to a tempting extrapolation made in the present investigation,
not source errata.  They do not refute the optimal-prefix conjecture:
both displayed determinants are nonzero.  A further proof for the
remaining residue class must therefore allow different valuation
behavior rather than extend the formula beyond its proved scope.

% Self-contained research section; standard amsmath/amsthm packages suffice.
% Bibliography keys and source notes are supplied in centered_cubic_notes.txt.
% No claim of historical priority is made for the lowest cubic moment.

\section{All regular odd centered cubic trigamma moments}
\label{cc:sec}

The canonical centered-square chapter evaluates all regular nonpositive
even moments of $\psi'(n+1)^2$.  Its cubic counterpart has the opposite
parity and a slightly larger, but still fixed, coordinate space.  We give
an explicit formula for every regular negative odd value, ordinary
absolutely convergent identities realizing these values, and a finite
polylogarithmic expression for their exponential generating function.

The cubic series itself has classical antecedents.  In particular,
Furdui evaluates $\sum_{n\ge1}n\psi'(n)^3$ in
\cite[Theorem 1(b)]{cc:furdui2015}; Hoffman's sums-of-tails calculus
provides a general multiple-zeta framework \cite{cc:hoffman2016}.
The contribution here is the entire centered polynomial-subtracted
family, its finite Bernoulli formula, and its explicit generator.
The lowest convergent evaluation is included as its normalization,
without a historical-priority claim.

\subsection{Regular values and ordinary convergent sums}

Put
\begin{equation}\label{cc:definitions}
 A_n=\psi'(n+1)=\zeta(2)-H_n^{(2)},\qquad
 x_n=n+\tfrac12,\qquad
 \mathcal U(s)=\sum_{n=0}^{\infty}A_n^3x_n^{-s}.
\end{equation}
The last series converges normally in $\Re s>-2$.  Define rational
numbers
\begin{equation}\label{cc:g}
 g_k=\sum_{\substack{a+b+c=k\\a,b,c\ge0}}
 B_{2a}(1/2)B_{2b}(1/2)B_{2c}(1/2),\qquad k\ge0.
\end{equation}
Thus $g_0=1$, $g_1=-1/4$, and $g_2=13/120$.
The Bernoulli convention throughout is $B_1=-1/2$.

\begin{proposition}[Continuation and the correct parity]
\label{cc:continuation}
The function $\mathcal U$ has a meromorphic continuation to the complex
plane, with simple poles at $s=-2-2k$, $k\ge0$, and residues $g_k$.
Every $s=-(2r+1)$, $r\ge0$, is regular.  Writing
\begin{equation}\label{cc:Mr}
 U_r=\mathcal U(-2r-1),
\end{equation}
one has the ordinary convergent identity
\begin{equation}\label{cc:ordinary}
 U_r=\sum_{n=0}^{\infty}
 \left\{x_n^{2r+1}A_n^3-
       \sum_{k=0}^{r-1}g_kx_n^{2r-2-2k}\right\}.
\end{equation}
The inner sum is empty for $r=0$; every bracket is $O(x_n^{-2})$.
\end{proposition}

\begin{proof}
The centered Hurwitz expansion gives, to arbitrary finite order,
\[
 A_n\sim\sum_{j\ge0}B_{2j}(1/2)x_n^{-2j-1},\qquad
 A_n^3\sim\sum_{k\ge0}g_kx_n^{-2k-3}.
\]
Subtracting finitely many terms and adding their Hurwitz-zeta sums
continues $\mathcal U$ to successive left half-planes.  The only
possible poles are the stated simple poles, with the stated residues.
For $j\ge1$, $B_{2j}(1/2)$ has sign $(-1)^j$, so every summand
in $g_k$ has sign $(-1)^k$; thus $g_k\ne0$.
At $s=-2r-1$ the subtracted sums are
$g_k\zeta(2k+2-2r,1/2)=0$ for $0\le k<r$.
Indeed $\zeta(-2m,1/2)=0$ for all $m\ge0$.
This proves \eqref{cc:ordinary} and its convergence.
\end{proof}

\subsection{A finite formula in four constant coordinates}

Set
\begin{align}
 C_1&=9\zeta(3)-\frac{17}{4}\zeta(4),
       \label{cc:C1}\\
 R_j&=-\frac{3}{2(2j-1)}
       \sum_{h=0}^{j-2}(2h+1)\binom{2j-1}{2h+3}
          B_{2j-4-2h}B_{2h}
       +\frac14\delta_{j,3},\qquad j\ge2,
       \label{cc:Rj}\\
 C_j&=9B_{2j-2}\zeta(3)-\frac12\delta_{j,2}\zeta(2)+R_j,
       \qquad j\ge2.
       \label{cc:Cj}
\end{align}
Here $\delta$ is the Kronecker delta.  The first few coordinates are
\begin{align}
 C_2&=\frac32\zeta(3)-\frac12\zeta(2)-\frac12,
       \notag\\
 C_3&=-\frac3{10}\zeta(3)-\frac25,
       \notag\\
 C_4&=\frac3{14}\zeta(3)-\frac5{56},
       \notag\\
 C_5&=-\frac3{10}\zeta(3)+\frac7{45}.
       \label{cc:Cexamples}
\end{align}

\begin{theorem}[All centered cubic moments]
\label{cc:all}
For every integer $r\ge0$,
\begin{equation}\label{cc:allformula}
 \boxed{\displaystyle
 U_r=\frac1{2r+2}\sum_{j=1}^{r+1}
      \binom{2r+2}{2j}B_{2r+2-2j}(1/2)C_j.}
\end{equation}
In particular every ordinary sum in \eqref{cc:ordinary} belongs to
\[
 \mathbb Q+\mathbb Q\zeta(2)+\mathbb Q\zeta(3)+\mathbb Q\zeta(4).
\]
No multiple-zeta coordinate of increasing weight is required as $r$
increases.
\end{theorem}

For example,
\begin{align}
 U_0&=\frac92\zeta(3)-\frac{17}{8}\zeta(4),
       \label{cc:U0}\\
 U_1&=-\frac18\zeta(2)-\frac34\zeta(3)
          +\frac{17}{32}\zeta(4)-\frac18,
       \label{cc:U1}\\
 U_2&=\frac5{48}\zeta(2)+\frac{47}{160}\zeta(3)
          -\frac{119}{384}\zeta(4)+\frac3{80},
       \label{cc:U2}\\
 U_3&=-\frac{49}{384}\zeta(2)-\frac{257}{1120}\zeta(3)
          +\frac{527}{1536}\zeta(4)-\frac{99}{4480}.
       \label{cc:U3}
\end{align}
Thus \eqref{cc:U1} evaluates
$\sum_{n\ge0}(x_n^3\psi'(n+1)^3-1)$, and
\eqref{cc:U2} evaluates
$\sum_{n\ge0}(x_n^5\psi'(n+1)^3-x_n^2+1/4)$.
These are convergent series; the displayed subtractions are part of
the identities.

\subsection{Proof by finite summation and explicit cancellation}

Let $\operatorname{CT}_N$ denote the coefficient of
$N^0(\log N)^0$ in a cutoff asymptotic expansion.  We use it only
after finite identities and standard asymptotic expansions.
Every partial sum of a nonnegative even power of $x_n$ has zero
constant coefficient as a polynomial in $N$.  Consequently
\eqref{cc:ordinary} implies
\begin{equation}\label{cc:cutoff}
 U_r=\operatorname{CT}_N\sum_{n=0}^{N-1}x_n^{2r+1}A_n^3.
\end{equation}
The centered power sum is the even polynomial
\begin{equation}\label{cc:centeredsum}
 \sum_{n=0}^{k-1}x_n^{2r+1}
 =\frac{B_{2r+2}(k+1/2)-B_{2r+2}(1/2)}{2r+2}
 =\frac1{2r+2}\sum_{j=1}^{r+1}
      \binom{2r+2}{2j}B_{2r+2-2j}(1/2)k^{2j}.
\end{equation}
Since $A_{k-1}=A_k+k^{-2}$, finite summation by parts reduces
\eqref{cc:allformula} to
\begin{equation}\label{cc:Ccutoff}
 C_j=\operatorname{CT}_N\left\{
 3\sum_{k=1}^Nk^{2j-2}A_k^2
 +3\sum_{k=1}^Nk^{2j-4}A_k
 +\sum_{k=1}^Nk^{2j-6}+N^{2j}A_N^3\right\}.
\end{equation}
It remains to prove that this expression is
\eqref{cc:C1}--\eqref{cc:Cj}.

\paragraph{The first two coordinates.}
The elementary tail rearrangements
\[
 \sum_{k=1}^{\infty}A_k^2=3\zeta(3)-\frac52\zeta(4),\qquad
 \sum_{k=1}^{\infty}\frac{A_k}{k^2}=\zeta_{\mathrm{MZ}}(2,2)=\frac34\zeta(4)
\]
give \eqref{cc:C1}.  For the first equality, include $k=0$ and
partition the two tail indices to obtain
$2\zeta_{\mathrm{MZ}}(2,1)+\zeta(3)=3\zeta(3)$, then remove $A_0^2$.

For $j=2$, write
$U_m(N)=\sum_{k=1}^Nk^mA_k$ and
$V_m(N)=\sum_{k=1}^Nk^mA_k^2$.
The same finite summation identities give
\begin{align*}
 \operatorname{CT}_N U_0(N)&=\gamma-\zeta(2)+1,\\
 \operatorname{CT}_N V_2(N)&=-\gamma+\frac12\zeta(2)
                         +\frac12\zeta(3)-\frac23,\\
 \operatorname{CT}_N(N^4A_N^3)&=-\frac32.
\end{align*}
For clarity, the first line follows from
$U_0(N)=H_N-H_N^{(2)}+NA_N$.
The second follows by using
$\sum_{k<\ell}k^2=\ell^3/3-\ell^2/2+\ell/6$ in the
quadratic summation identity below; the needed additional constants are
$\operatorname{CT}_NU_1(N)=1/4-\gamma/2$ and
$U_{-1}(\infty)=\zeta(3)$.
Substitution in \eqref{cc:Ccutoff} yields the displayed $C_2$.

\paragraph{The general coordinate.}
Let $m=2j-2\ge4$ and put
\[
 F_m(N)=\sum_{k=1}^Nk^m
 =\sum_{\substack{a=1\\a\text{ odd}}}^{m+1}f_aN^a+\frac12N^m,
 \qquad
 f_a=\frac1{m+1}\binom{m+1}{a}B_{m+1-a}.
\]
Then
\[
 F_m(k-1)=\sum_{\substack{a=1\\a\text{ odd}}}^{m+1}f_ak^a-\frac12k^m.
\]
Finite summation by parts gives the exact relation
\begin{equation}\label{cc:quadraticcutoff}
 V_m(N)=F_m(N)A_N^2+
 2\sum_{k=1}^N\frac{F_m(k-1)A_k}{k^2}
 +\sum_{k=1}^N\frac{F_m(k-1)}{k^4}.
\end{equation}
Its term $-U_{m-2}(N)$ cancels exactly against the corresponding
linear-tail term in \eqref{cc:Ccutoff}.  This cancellation occurs
before constants or divergent harmonic terms are taken.

For every positive odd $q$, another finite summation gives
\begin{equation}\label{cc:oddlinear}
 \operatorname{CT}_NU_q(N)
   =\frac{2-q}{4}B_{q-1}-\frac12\delta_{q,1}\gamma.
\end{equation}
Indeed, if $F_q(N)=\sum_{k=1}^Nk^q$, then
\[
 U_q(N)=F_q(N)A_N+\sum_{k=1}^N\frac{F_q(k-1)}{k^2}.
\]
Use the uncentered expansion
$A_N\sim\sum_{a\ge1}B_{a-1}N^{-a}$.
Only the $N^2$ and $N^q$ terms of $F_q$ can contribute to the
constant when $q\ge3$, yielding $(2-q)B_{q-1}/4$.
The case $q=1$ follows directly from
$F_1(N)=(N^2+N)/2$ and $(N-H_N)/2$.

In \eqref{cc:quadraticcutoff}, the $a=1$ terms contribute
$9f_1\zeta(3)=9B_m\zeta(3)$ to \eqref{cc:Ccutoff}.
The $a=3$ harmonic constants cancel by \eqref{cc:oddlinear}; all
remaining nonnegative-power sums have constant zero.
For odd $a\ge3$,
\[
 [N^{-a}]A_N^2=-B_{a-3}.
\]
Combining these boundary coefficients with \eqref{cc:oddlinear}
leaves the rational sum
\[
 \frac32\sum_{\substack{a=3\\a\text{ odd}}}^{m+1}
          f_a(2-a)B_{a-3},
\]
apart from the even boundary coefficient.  The latter is exactly
\begin{equation}\label{cc:boundaryexception}
 \frac32[N^{-m}]A_N^2+[N^{-m-2}]A_N^3
       =\frac14\delta_{m,4}.
\end{equation}
To see this, all odd-index Bernoulli numbers except $B_1=-1/2$
vanish.  The terms containing exactly one $B_1$ cancel the ordinary
even-even convolution; the only surviving possibility is the
$B_1^2/B_1^3$ contribution at $m=4$.
Substitute $a=2h+3$ and $m=2j-2$ to obtain
\eqref{cc:Rj}--\eqref{cc:Cj}.  This proves
Theorem~\ref{cc:all}.

\subsection{A fixed coordinate space and a rational cubic identity}

Treat $\zeta(2),\zeta(3),\zeta(4),1$ as formal coordinates.
The coefficient matrix of $U_0,U_1,U_2,U_3$ in this order is
\begin{equation}\label{cc:matrix}
 M=\begin{pmatrix}
 0&9/2&-17/8&0\\
 -1/8&-3/4&17/32&-1/8\\
 5/48&47/160&-119/384&3/80\\
 -49/384&-257/1120&527/1536&-99/4480
 \end{pmatrix},\qquad
 \det M=\frac{51}{81920}.
\end{equation}
Thus these first four values generate the entire formal coordinate
space of the cubic family, and every later value is an explicit
rational combination of them.  This is a coefficient-space statement;
it asserts no unproved arithmetic independence of the four constants.

One particularly simple elimination is
\begin{equation}\label{cc:rationalmoment}
 11U_0+198U_1+512U_2+224U_3=-\frac{21}{2}.
\end{equation}
Consequently, with
$P(x)=11x+198x^3+512x^5+224x^7$, one obtains the entirely rational
ordinary sum
\begin{equation}\label{cc:rationalsum}
 \boxed{\displaystyle
 \sum_{n=0}^{\infty}\left\{
 P(x_n)\psi'(n+1)^3-224x_n^4-456x_n^2-\frac{1414}{15}
 \right\}=-\frac{21}{2}.}
\end{equation}
Its summand is $O(x_n^{-2})$.  Equation \eqref{cc:rationalsum}
follows from the four proved values and their prescribed subtractions,
so no numerical relation search or independence assumption enters it.

Another elimination isolates the odd zeta coordinate:
\begin{equation}\label{cc:zeta3moment}
37U_0+2136U_1+11344U_2+7168U_3=252\zeta(3).
\end{equation}
With $P_3(x)=37x+2136x^3+11344x^5+7168x^7$, this becomes
\begin{equation}\label{cc:zeta3sum}
 \boxed{\begin{aligned}
 \sum_{n=0}^{\infty}\Bigl\{P_3(x_n)\psi'(n+1)^3
  -7168x_n^4-9552x_n^2-\frac{1148}{15}\Bigr\}
  =252\zeta(3).
 \end{aligned}}
\end{equation}
It is again an ordinary absolutely convergent sum. Equations
\eqref{cc:rationalsum} and \eqref{cc:zeta3sum} demonstrate that
the fixed four-coordinate description supplies exact extraction
identities, in addition to individual moment evaluations.

\subsection{The cubic generator from the known square generator}

Let
\[
 \mathcal G_2(t)=\sum_{r\ge0}\mathcal T(-2r)\frac{t^{2r}}{(2r)!},
 \qquad
 \mathcal G_3(t)=\sum_{r\ge0}U_r\frac{t^{2r+1}}{(2r+1)!},
\]
where $\mathcal T(s)=\sum_{n\ge0}A_n^2x_n^{-s}$ for $\Re s>-1$,
continued meromorphically elsewhere, is the centered-square Dirichlet
series of the canonical chapter. The inherited explicit expression is
\begin{equation}\label{cc:G2}
 \mathcal G_2(t)=\frac{t}{2\sinh(t/2)}
 \left[3\zeta(3)-\frac1{4t}
       \int_0^t(t-u)u^2\coth(u/2)\,du\right].
\end{equation}
This is recalled with attribution to \cite{cc:repo-square}, not
presented as an additional new square-moment theorem.

\begin{theorem}[An analytic cubic moment generator]
\label{cc:generator}
The series $\mathcal G_3$ converges for $|t|<2\pi$.  On that disk,
\begin{equation}\label{cc:G3integral}
 \boxed{\displaystyle
 \mathcal G_3(t)=\frac1{2\sinh(t/2)}
 \int_0^t(t-u)\left\{
    3\cosh(u/2)\mathcal G_2(u)-\frac{17}{4}\zeta(4)
    -\frac14\zeta(2)u^2+\frac{u^4}{96}
 \right\}\,du.}
\end{equation}
The expression is holomorphic and odd; its apparent singularity at
zero is removable.  The integration path lies in the disk and is
otherwise arbitrary.
\end{theorem}

\begin{proof}
Put $C(t)=\sum_{j\ge1}C_jt^{2j}/(2j)!$.
The Bernoulli transform \eqref{cc:allformula} gives
\[
 \mathcal G_3(t)=\frac{C(t)}{2\sinh(t/2)}.
\]
Write $E(t)=\tfrac t2\coth(t/2)$ and
\[
 F(t)=\frac1t\int_0^t(t-u)uE(u)\,du
     =\sum_{h\ge0}\frac{(2h+1)B_{2h}}{(2h+3)!}t^{2h+2}.
\]
Direct coefficient multiplication of \eqref{cc:Rj} yields
\begin{align*}
 C''(t)
 &=9\zeta(3)E(t)-\frac32E(t)F(t)
          -\frac{17}{4}\zeta(4)-\frac14\zeta(2)t^2+\frac{t^4}{96}\\
 &=3\cosh(t/2)\mathcal G_2(t)
          -\frac{17}{4}\zeta(4)-\frac14\zeta(2)t^2+\frac{t^4}{96}.
\end{align*}
In the second line use \eqref{cc:G2} and
$E(t)=\cosh(t/2)t/(2\sinh(t/2))$.
All these series converge in $|t|<2\pi$ by the Bernoulli generating
functions.  The initial conditions are $C(0)=C'(0)=0$.
Twice integrating proves \eqref{cc:G3integral} and the stated
analyticity and parity.
\end{proof}

\subsection{A finite depth-two polylogarithmic expression}

For $0<z<1$ use the convention
\[
 \operatorname{Li}_{a,b}(z)
   =\sum_{m>n\ge1}\frac{z^m}{m^a n^b}.
\]
This is the one-variable multiple polylogarithm with the outer index
carrying the argument.  In the following formula, every displayed
depth-two first index is at least two.

\begin{theorem}[Closed polylogarithmic generator]
\label{cc:polylog}
For real $0<t<2\pi$, set $z=e^{-t}$ and define
\begin{align}
 \mathcal C(t)={}&
 \left(9\zeta(4)+6t\zeta(3)-\frac{t^4}{16}\right)
           \operatorname{Li}_2(z)
 +\left(12\zeta(3)-\frac{t^3}{2}\right)\operatorname{Li}_3(z)
       \notag\\
 &-3t^2\operatorname{Li}_4(z)-12t\operatorname{Li}_5(z)
       -\frac{45}{2}\operatorname{Li}_6(z)
       -\frac32t^2\operatorname{Li}_{2,2}(z)
       \notag\\
 &-6t\left(\operatorname{Li}_{3,2}(z)+\operatorname{Li}_{2,3}(z)\right)
       -9\left(\operatorname{Li}_{4,2}(z)+\operatorname{Li}_{2,4}(z)\right)
       -12\operatorname{Li}_{3,3}(z)
       \notag\\
 &+\frac{15}{2}\zeta(6)-6\zeta(3)^2+6t\zeta(5)
       +\frac12t^3\zeta(3)+\frac18t^2\zeta(4)
       -\frac1{48}t^4\zeta(2)-\frac{t^6}{1440}.
       \label{cc:Cpolylog}
\end{align}
Then
\begin{equation}\label{cc:G3polylog}
 \mathcal G_3(t)=\frac{\mathcal C(t)}{2\sinh(t/2)}.
\end{equation}
Equation \eqref{cc:G3integral} specifies its single-valued holomorphic
odd continuation throughout $|t|<2\pi$.
\end{theorem}

\begin{proof}
For ordinary polylogarithms,
$\frac{d}{dt}\operatorname{Li}_a(e^{-t})=-\operatorname{Li}_{a-1}(e^{-t})$.
For the depth-two convention just stated,
\[
 \frac{d}{dt}\operatorname{Li}_{a,b}(e^{-t})
   =-\operatorname{Li}_{a-1,b}(e^{-t})\quad(a\ge2),\qquad
 \frac{d}{dt}\operatorname{Li}_{1,b}(e^{-t})
   =-\frac{\operatorname{Li}_b(e^{-t})}{e^t-1}.
\]
Together with $\coth(t/2)=(1+e^{-t})/(1-e^{-t})$, these give,
by finite differentiation, that $\mathcal C''$ is the integrand in
\eqref{cc:G3integral}.

The integration constants may be checked by anchored primitives.
If $F_a'(t)=-F_{a-1}(t)$ and the displayed endpoint values exist, then
\begin{equation}\label{cc:anchoredprimitive}
 \int_0^t(t-u)u^kF_0(u)\,du
 =\sum_{j=0}^k(j+1)\frac{k!}{(k-j)!}t^{k-j}F_{j+2}(t)
       -(k+1)!F_{k+2}(0)+tk!F_{k+1}(0).
\end{equation}
Apply this with $F_a=\operatorname{Li}_{s+a}(e^{-t})$, or
$F_a=\operatorname{Li}_{a,b}(e^{-t})$ for $k\ge1$.
The only potentially divergent endpoint is kept in the convergent
combination
\[
 \operatorname{Li}_{1,4}(e^{-t})-\zeta(4)\operatorname{Li}_1(e^{-t})
 \longrightarrow -\zeta_{\mathrm{MZ}}(4,1)-\zeta(5)
                 =\zeta(2)\zeta(3)-3\zeta(5).
\]
The elementary depth-two reductions
\begin{align*}
 \zeta_{\mathrm{MZ}}(2,3)&=\tfrac92\zeta(5)-2\zeta(2)\zeta(3),&
 \zeta_{\mathrm{MZ}}(3,2)&=3\zeta(2)\zeta(3)-\tfrac{11}{2}\zeta(5),\\
 \zeta_{\mathrm{MZ}}(4,2)+\zeta_{\mathrm{MZ}}(2,4)&=\zeta(2)\zeta(4)-\zeta(6),&
 2\zeta_{\mathrm{MZ}}(3,3)&=\zeta(3)^2-\zeta(6)
\end{align*}
reduce the endpoint constants to those in \eqref{cc:Cpolylog}.
Thus $\mathcal C(0)=\mathcal C'(0)=0$ in the limiting sense,
and \eqref{cc:G3integral} proves the claim.  No differentiation of
an identity known only at integral zeta orders is used.
\end{proof}

\subsection{Verification scope and immediate continuation questions}

The supplied script independently reconstructs the cutoff constants
using Faulhaber polynomials and the uncentered trigamma expansion; it
does not insert \eqref{cc:Rj} into that reconstruction.  It agrees
with the closed formula for $1\le j\le18$, checks the generator
coefficient conversion over the same range, and verifies the second
derivative of \eqref{cc:Cpolylog} as an exact symbolic zero.
Independent direct tail sums plus centered Hurwitz-tail corrections
check $0\le r\le8$.  The largest displayed residual in the higher
cutoff calculation is below $5\cdot10^{-76}$ at 110-digit working
precision.  Those residuals are numerical diagnostics, not certified
interval bounds or replacements for the finite-sum proof.

The new result suggests three concrete questions.  First, determine
whether all regular centered moments of $\psi'(n+1)^4$ admit a
similarly explicit fixed coordinate space, and derive its generator
with every subtraction and endpoint constant specified.  Second,
identify the analogous regular-parity moment spaces for mixed products
$\prod_i\psi^{(m_i)}(n+1)$, distinguishing formal coefficient rank
from arithmetic independence.  Third, determine whether the particular
depth-two combination in \eqref{cc:Cpolylog} reduces to a prescribed
finite class of classical polylogarithms; a failed numerical search
would not prove that it cannot.  None of these questions is settled
by the cubic calculation above.

% Standalone section. Requires amsmath, amssymb, amsthm and hyperref.
% The parent article should provide theorem, proposition, corollary, remark.
% Bibliography keys: dhr20261010, dlmfgauss, dlmfgammaratio, dlmfhurwitz, dlmfpolylog.
\section{The next two-particle resonance: overlap subtraction and exact constants}
\label{sec:second-resonance}

The determinantal harmonic-resonance report explicitly proposed the next
two-particle resonance as the first test of a higher-pole subtraction
\cite[Section~9]{incoming-determinantal}. We resolve that target for arbitrary positive
shape parameters, and by joint holomorphy for their right half-planes.
The result consists of an explicit normally convergent subtraction, its
complete spectral derivative calculus, and a finite-cutoff constant proved
independently from an exact discrete Pearson identity. These statements
are proved here; no claim of exhaustive literature novelty is made.

\subsection{Definitions and a finite identity}
Let $\Re a,\Re c>0$ and put
\begin{align}
 w_n(u)&=\frac{\Gamma(n+a)\Gamma(n+c)}
 {\Gamma(n+a+c+1+u)\Gamma(n+1)},\label{sr:eq:weight}\\
 S_N(u)&=\sum_{0\le m<n\le N}(n-m)^2w_m(u)w_n(u),\qquad
 K_n(u)=\sum_{0\le m<n}(n-m)^2w_m(u)w_n(u).
 \label{sr:eq:cutoff}
\end{align}
Thus $S_N=\sum_{n=0}^N K_n$ and $K_0=0$. There is no factor $2!$.
Use the abbreviations
\begin{gather}
 \mu=\frac1{ac},\qquad h=\frac1a+\frac1c,\qquad q=1+h,
 \qquad s=ac+a+c,\label{sr:eq:abbreviations}\\
 \rho=\psi(1)-\psi(a+1)-\psi(c+1),\qquad
 \tau=\psi'(1)-\psi'(a+1)-\psi'(c+1),\qquad r_2=\rho^2+\tau.
 \label{sr:eq:rho}
\end{gather}
The letter $q$ in this section is a scalar abbreviation, not a
cyclotomic level or an integration order.

\begin{proposition}[Exact Pearson reduction]\label{sr:prop:pearson}
Write $M_j(N,u)=\sum_{n=0}^N n^jw_n(u)$ and
$t_n(u)=n(n+a+c+u)w_n(u)$. For every $N\ge0$,
\begin{align}
 acM_0(N,u)-uM_1(N,u)&=t_{N+1}(u),\label{sr:eq:pearson0}\\
 (1-u)M_2(N,u)+(s+u)M_1(N,u)&=Nt_{N+1}(u).
 \label{sr:eq:pearson1}
\end{align}
Consequently, at $u=0$,
\begin{align}
 acM_0(N,0)&=(N+1)(N+1+a+c)w_{N+1}(0),\label{sr:eq:M0exact}\\
 S_N(0)&=acN M_0(N,0)^2-sM_0(N,0)M_1(N,0)-M_1(N,0)^2.
 \label{sr:eq:finite-det}
\end{align}
\end{proposition}
\begin{proof}
The exact weight ratio is
\[
 \frac{w_{n+1}(u)}{w_n(u)}
 =\frac{(n+a)(n+c)}{(n+a+c+1+u)(n+1)}.
\]
It gives
\[
 t_{n+1}-t_n=(ac-un)w_n,
 \qquad
 (n+1)t_{n+1}-nt_n=((1-u)n^2+sn+ac)w_n.
\]
Sum these identities, use $t_0=0$, and eliminate $acM_0$ to obtain
\eqref{sr:eq:pearson0}--\eqref{sr:eq:pearson1}. The elementary identity
$S_N=M_0M_2-M_1^2$ follows by expanding $(n-m)^2$ and pairing the two
strict orderings. Substitution at zero proves \eqref{sr:eq:finite-det}.
\end{proof}

\subsection{The Gamma product and its two different singular scales}
Define holomorphic germs
\begin{align}
 R(u)&=\frac{\Gamma(a+1)\Gamma(c+1)\Gamma(1+u)}
 {\Gamma(a+1+u)\Gamma(c+1+u)},& A(u)&=\mu R(u),\label{sr:eq:R}\\
 \beta(u)&=-s-(a+c+\tfrac12)u-\tfrac12u^2,
 &V(u)&=\frac1{1-u^2},\label{sr:eq:beta}\\
 P(u)&=(1+hu+\mu u^2)R(u)^2,
 &T(u;\alpha)&=A(u)(\beta(u)+\alpha u).
 \label{sr:eq:PT}
\end{align}
In particular $R(0)=1$, $R'(0)=\rho$, and $R''(0)=r_2$.
For $\Re u>1$, the partition converges absolutely, and Gauss summation
\cite[15.4.20]{dlmfgauss} gives
\begin{equation}\label{sr:eq:partition-product}
 \mathcal Z(u):=\sum_{n\ge0}K_n(u)
 =\frac{P(u)}{u^2(u-1)}.
\end{equation}
For completeness, its three moment sums in this domain are
\[
 M_0(\infty,u)=\mu R(u),\qquad
 M_1(\infty,u)=\frac{R(u)}u,\qquad
 M_2(\infty,u)=\frac{R(u)(s+u)}{u(u-1)}.
\]
Their determinant is \eqref{sr:eq:partition-product}. Thus the product
needs only ordinary Gauss summation, not an assumed higher-dimensional
Selberg formula.

Fix $\Re\alpha>0$ and write $y=n+\alpha$, with the logarithm on the
right half-plane. Define the \emph{single grouped counterterm}
\begin{align}
 Q_n(u;\alpha)={}&A(u)y^{-u}+T(u;\alpha)y^{-1-u}\notag\\
 &+\frac2u\left\{V(u)y^{-1-2u}-R(u)y^{-1-u}\right\}.
 \label{sr:eq:Q}
\end{align}
The last line is holomorphic at $u=0$ because its numerator vanishes
there. The exponent $1+2u$ is essential: it accounts for the region in
which the inner summation index also becomes large. Replacing it by
$1+u$ loses the logarithmic overlap term.

\begin{theorem}[Normally convergent second-resonance germ]
\label{sr:thm:germ}
For $\Re a,\Re c,\Re\alpha>0$, the series
\begin{equation}\label{sr:eq:D}
 \mathcal D(u;\alpha)=\sum_{n\ge0}\{K_n(u)-Q_n(u;\alpha)\}
\end{equation}
converges absolutely and locally normally for $u$ in a neighborhood of
zero. It is jointly holomorphic in these parameters, and every fixed
mixed parameter derivative can be taken termwise. Its exact value is
\begin{align}
 \mathcal D(u;\alpha)={}&\frac{P(u)}{u^2(u-1)}-A(u)\zeta(u,\alpha)
 -T(u;\alpha)\zeta(1+u,\alpha)\notag\\
 &+\frac{2R(u)}u\zeta(1+u,\alpha)
 -\frac{2V(u)}u\zeta(1+2u,\alpha).
 \label{sr:eq:master}
\end{align}
All apparent poles at $u=0$ cancel in this complete expression.
\end{theorem}
\begin{proof}
Work on compact shape-parameter sets and choose $0<\delta<1/4$.
The locally uniform Gamma-ratio expansion \cite[5.11.13]{dlmfgammaratio}
is
\begin{equation}\label{sr:eq:w-asymptotic}
 w_n(u)=n^{-2-u}\left(1+\frac{\beta(u)}n+O(n^{-2})\right),
 \qquad |u|\le\delta.
\end{equation}
Gauss summation gives $\sum_{m\ge0}w_m(u)=A(u)$ already in
$\Re u>-1$. Hence, uniformly on the chosen compact sets,
\begin{equation}\label{sr:eq:lower-M0}
 \sum_{m<n}w_m(u)=A(u)-\frac{n^{-1-u}}{1+u}+O(n^{-2+\delta}).
\end{equation}
For the first moment, the series
$\sum_{m\ge1}\{mw_m(u)-m^{-1-u}\}$ converges locally normally for
$\Re u>-1$. On $\Re u>0$ its value is $R(u)/u-\zeta(1+u)$, so that
identity continues through zero. Its tail, together with the leading
Hurwitz-zeta tail, gives the \emph{holomorphic} estimate
\begin{equation}\label{sr:eq:lower-M1}
 \sum_{m<n}mw_m(u)=\frac{R(u)-n^{-u}}u+O(n^{-1+\delta}).
\end{equation}
The quotient has its removable value at zero. Finally, summing
$m^2w_m(u)=m^{-u}+O(m^{-1+\delta})$ gives
\begin{equation}\label{sr:eq:lower-M2}
 \sum_{m<n}m^2w_m(u)=\frac{n^{1-u}}{1-u}
 +O(n^\delta\log(n+1)).
\end{equation}
Insert these estimates in
\[
 K_n=w_n\left(n^2\sum_{m<n}w_m-2n\sum_{m<n}mw_m
                         +\sum_{m<n}m^2w_m\right).
\]
The coefficient of $n^{-1-2u}$ is
\[
 -\frac1{1+u}+\frac2u+\frac1{1-u}
 =\frac{2}{u(1-u^2)}.
\]
The resulting paired expression is holomorphic at zero, and
\[
 K_n=A(u)n^{-u}+A(u)\beta(u)n^{-1-u}
 +\frac2u\{V(u)n^{-1-2u}-R(u)n^{-1-u}\}
 +O(n^{-2+2\delta}\log(n+1)).
\]
Changing $n$ to $y=n+\alpha$ in the leading scale adds
$\alpha uA(u)y^{-1-u}$; the remaining changes are summable uniformly.
This proves \eqref{sr:eq:D}, since the displayed error is summable.
Cauchy's integral formula on slightly larger compact parameter sets
proves the assertions for every fixed derivative.

To identify the sum, the same estimates hold on compact subsets of
$-\delta<\Re u<1+\delta$, excluding $u=1$. In particular,
\eqref{sr:eq:lower-M2} still holds there: summing $m^{-u}$ gives
$n^{1-u}/(1-u)$ plus a locally bounded constant and
$O(n^\delta)$, while the partial-sum estimate for
$m^2w_m-m^{-u}$ is $O(n^\delta\log(n+1))$. Constants are uniform
on compact sets avoiding $u=1$. Thus the subtracted remainder remains
summable throughout this strip. Both complete expressions are
holomorphic on its connected punctured domain, including their
removable values at zero.
On $1<\Re u<1+\delta$ all original sums and counterterms converge
separately, and their difference is \eqref{sr:eq:master} by
\eqref{sr:eq:partition-product}. The identity theorem gives the
formula near zero. Its removability there was established by the
normally convergent left side, independently of formal Laurent algebra.
\end{proof}

\subsection{Exact hard-cutoff constant and spectral comparison}
Set
\begin{equation}\label{sr:eq:lambda-C}
 \lambda=-q-2\rho,\qquad C=-\rho^2-q\rho-2.
\end{equation}
\begin{theorem}[Closed constant at the next resonance]
\label{sr:thm:constant}
As $N\to\infty$,
\begin{equation}\label{sr:eq:cutoff-asymptotic}
 \boxed{\displaystyle
 S_N(0)=\mu N-\log^2N+\lambda\log N+C+o(1).}
\end{equation}
The error is $O(\log N/N)$ for fixed shape parameters, locally uniformly
on their right half-planes. The spectral Laurent expansion is instead
\begin{equation}\label{sr:eq:spectral-laurent}
 \mathcal Z(u)=-\frac1{u^2}+\frac\lambda u+F+O(u),\qquad
 F=-\mu-q-2q\rho-2\rho^2-\tau.
\end{equation}
Thus the exact correction between the two constants is
\begin{equation}\label{sr:eq:finitepart-correction}
 F-C=-\mu-q+2-q\rho-\rho^2-\tau.
\end{equation}
\end{theorem}
\begin{proof}
At zero, \eqref{sr:eq:lower-M0}--\eqref{sr:eq:lower-M1} give
\[
 M_0(N,0)=\mu-\frac1N+O(N^{-2}),\qquad
 M_1(N,0)=\log N+\rho+O(N^{-1}).
\]
Substituting them in the exact finite identity
\eqref{sr:eq:finite-det} gives
\[
 acNM_0^2=\mu N-2+O(N^{-1}),\quad
 sM_0M_1=q(\log N+\rho)+O(\log N/N),
\]
and the asserted cutoff formula follows. This proof does not identify
a spectral finite part with a cutoff constant. Separately expanding
$R(u)=1+\rho u+(\rho^2+\tau)u^2/2+O(u^3)$ in
\eqref{sr:eq:partition-product} proves \eqref{sr:eq:spectral-laurent}
and \eqref{sr:eq:finitepart-correction}.
\end{proof}

With the Stieltjes convention
\[
 \zeta(1+u,\alpha)=\frac1u+
 \sum_{k\ge0}\frac{(-1)^k}{k!}\gamma_k(\alpha)u^k,
 \qquad \gamma_0(\alpha)=-\psi(\alpha),
\]
the constant term of Theorem~\ref{sr:thm:germ} is the ordinary
single-bracket identity
\begin{equation}\label{sr:eq:D0}
 \boxed{\displaystyle
 \sum_{n\ge0}\left\{K_n(0)-\mu-
 \frac{\lambda-2\log(n+\alpha)}{n+\alpha}\right\}
 =C-\mu-\lambda\gamma_0(\alpha)+2\gamma_1(\alpha).}
\end{equation}
Its agreement with \eqref{sr:eq:cutoff-asymptotic} also follows by
summing the elementary counterterms before taking the limit. The
coefficient $2\gamma_1$ comes from the doubled exponent $2u$.

\subsection{Every derivative and an explicit first row}
Write $R_j=R^{(j)}(0)$, $T_j=\partial_u^jT(0;\alpha)$ and
$V_j=V^{(j)}(0)$, so $V_j=j!$ for even $j$ and $V_j=0$ for odd $j$.
Let
\begin{equation}\label{sr:eq:N-germ}
 \mathcal N(u)=\frac{P(u)}{u-1}+2R(u)-V(u)-uT(u;\alpha).
\end{equation}
Direct substitution gives $\mathcal N(0)=\mathcal N'(0)=0$.
Thus $\mathcal N(u)/u^2$ is an explicitly given holomorphic
Gamma/polygamma germ, not an additional unspecified constant.

\begin{corollary}[Finite formula for every spectral row]
\label{sr:cor:alljets}
For every $k\ge0$, put $D_k(\alpha)=\partial_u^k\mathcal D(0;\alpha)$.
Then
\begin{align}
 D_k(\alpha)={}&
 \frac{\mathcal N^{(k+2)}(0)}{(k+1)(k+2)}
 -\mu\sum_{j=0}^k\binom kj R_j\zeta^{(k-j)}(0,\alpha)\notag\\
 &-\sum_{j=0}^k\binom kj T_j(-1)^{k-j}\gamma_{k-j}(\alpha)\notag\\
 &+\frac2{k+1}\sum_{j=0}^{k+1}\binom{k+1}j
 \bigl(R_j-2^{k+1-j}V_j\bigr)(-1)^{k+1-j}
 \gamma_{k+1-j}(\alpha).
 \label{sr:eq:alljets}
\end{align}
Here $\zeta^{(j)}$ denotes differentiation in the order variable.
The normally convergent left side is
$\sum_{n\ge0}\{K_n^{(k)}(0)-Q_n^{(k)}(0;\alpha)\}$. If
$L_n=\log(n+\alpha)$, every counterterm is the finite polynomial
\begin{align}
 Q_n^{(k)}(0;\alpha)={}&\mu\sum_{j=0}^k\binom kjR_j(-L_n)^{k-j}
 +\frac1{n+\alpha}\sum_{j=0}^k\binom kjT_j(-L_n)^{k-j}\notag\\
 &+\frac2{(k+1)(n+\alpha)}
 \sum_{j=0}^{k+1}\binom{k+1}j
 \left\{V_j(-2L_n)^{k+1-j}-R_j(-L_n)^{k+1-j}\right\}.
 \label{sr:eq:counterterm-jets}
\end{align}
\end{corollary}
\begin{proof}
Put $J(u;\alpha)=\zeta(1+u,\alpha)-1/u$. Rearranging
\eqref{sr:eq:master} gives the manifestly regular expression
\[
 \mathcal D=\frac{\mathcal N(u)}{u^2}-\mu R(u)\zeta(u,\alpha)
 -T(u;\alpha)J(u;\alpha)
 +\frac2u\{R(u)J(u;\alpha)-V(u)J(2u;\alpha)\}.
\]
Apply the product rule. Division of a germ vanishing to order $m$
converts its $(k+m)$th derivative to its $k$th derivative divided by
$(k+1)\cdots(k+m)$; this gives the two denominators in
\eqref{sr:eq:alljets}. The same argument applied to \eqref{sr:eq:Q}
gives \eqref{sr:eq:counterterm-jets}. Normal convergence was already
proved in Theorem~\ref{sr:thm:germ}.
\end{proof}

All coefficients in this formula are finite polygamma polynomials.
For example, the complete exponential Bell polynomials give
\[
 R_j=\operatorname{Bell}_j(\ell_1,\ldots,\ell_j),\quad
 \ell_j=\psi^{(j-1)}(1)-\psi^{(j-1)}(a+1)-\psi^{(j-1)}(c+1).
\]
The other coefficients follow by multiplication by the explicit
polynomials and rational functions in \eqref{sr:eq:beta}--\eqref{sr:eq:PT}.
In fact the apparently required derivative $R_{k+2}$ always cancels.
To see this without coefficient cancellation, put
$W(u)=(R(u)-1)/u$ and $d=\mu(\alpha-a-c-\tfrac12)$. Then
\begin{equation}\label{sr:eq:lower-order-G}
 \frac{\mathcal N(u)}{u^2}=-W(u)^2-V(u)
 +R(u)\left\{\frac{(q+\mu u)W(u)+q+\mu}{u-1}
                    -d+\frac\mu2u\right\}.
\end{equation}
Its $k$th derivative uses only $R_0,\ldots,R_{k+1}$.
Thus the $k$th identity needs shape polygammas only through order $k$,
while its Stieltjes constants run through order $k+1$.
The summand itself has the equally explicit expression
\begin{equation}\label{sr:eq:K-bell}
 K_n^{(k)}(0)=\sum_{m<n}(n-m)^2w_m(0)w_n(0)
 \operatorname{Bell}_k(x_1(m,n),\ldots,x_k(m,n)),
\end{equation}
where
\[
 x_j(m,n)=-\psi^{(j-1)}(m+a+c+1)-\psi^{(j-1)}(n+a+c+1).
\]
At positive integer parameters these are ordinary generalized harmonic
numbers and zeta constants; at arbitrary shifts they are their
polygamma continuations.

The first derivative can be printed without Bell notation. Define
\begin{align}
 B(\alpha)&=r_2+q\rho+q-3+\mu(\tfrac12-\alpha),\label{sr:eq:B}\\
 C_1&=-\rho r_2-q\rho^2-\frac q2r_2
 -(2\mu+q+1)\rho-q-\frac\mu2.
 \label{sr:eq:C1}
\end{align}
Then
\begin{align}
 Q_n'(0;\alpha)&=\mu(\rho-L_n)
 +\frac{3L_n^2-\lambda L_n-B(\alpha)}{n+\alpha},
 \label{sr:eq:Q1}\\
 D_1(\alpha)&=C_1-\mu\zeta'(0,\alpha)
 +B(\alpha)\gamma_0(\alpha)+\lambda\gamma_1(\alpha)
 -3\gamma_2(\alpha).
 \label{sr:eq:D1}
\end{align}
Here $K_n'(0)$ is minus the defining double summand multiplied by
$\psi(m+a+c+1)+\psi(n+a+c+1)$. The term
$\zeta'(0,\alpha)=\log\Gamma(\alpha)-\tfrac12\log(2\pi)$
is classical \cite[25.11.18]{dlmf-hurwitz}. Thus this row is an explicit
ordinary convergent identity connecting a polygamma-weighted double
sum, log-Gamma, and the first three generalized Stieltjes constants.

For comparison, the corresponding unsubtracted finite cutoff has
the full first-derivative expansion
\begin{align}
 S_N'(0)={}&-\mu N\log N+\mu(\rho+1)N+\log^3N
 -\frac\lambda2\log^2N\notag\\
 &-(r_2+q\rho+q-3+\mu)\log N+C_1+\mu\rho+o(1).
 \label{sr:eq:cutoff-derivative}
\end{align}
Indeed take $\alpha=1$ in \eqref{sr:eq:Q1}, sum through $N$,
use $\sum_{j=1}^{N+1}\log j=\log\Gamma(N+2)$ and the defining
Stieltjes limits, and apply Stirling's formula. The constants
$\gamma_0,\gamma_1,\gamma_2$ cancel against \eqref{sr:eq:D1}; the
$\log(2\pi)$ terms cancel as well.

\subsection{A uniform bridge from spectral germs to finite cutoffs}
For any fixed finite derivative order, one can obtain the full cutoff
expansion without making new interchanges. If $Y=N+1+\alpha$, the
normally convergent tail in Theorem~\ref{sr:thm:germ} gives
\begin{align}
 S_N(u)={}&\mathcal Z(u)-A(u)\zeta(u,Y)
 -\left\{T(u;\alpha)-\frac{2R(u)}u\right\}\zeta(1+u,Y)\notag\\
 &-\frac{2V(u)}u\zeta(1+2u,Y)
 +O(N^{-1+2\delta}\log(N+1)).
 \label{sr:eq:uniform-cutoff}
\end{align}
The complete expression has removable poles at zero. The error is
locally uniform for $|u|\le\delta<1/4$, after shrinking a surrounding
disk if needed for derivatives. Thus every fixed Taylor coefficient
may be extracted, with a remainder tending to zero. The standard
large-$Y$ Hurwitz expansion then generates every $N\log^jN$ term,
every logarithmic counterterm, and its constant. The overlap exponent
$2u$ controls the higher logarithmic powers.

\subsection{Shift differentiation and an anchored primitive}
The subtraction point is auxiliary; it does not move the original
particle grid. Nevertheless \eqref{sr:eq:D0} has the exact laws
\begin{align}
 D_0(\alpha+1)-D_0(\alpha)
 &=\frac{\lambda-2\log\alpha}{\alpha},\label{sr:eq:alpha-shift}\\
 \frac{d}{d\alpha}D_0(\alpha)
 &=(\lambda+2)\zeta(2,\alpha)+2\zeta'(2,\alpha).
 \label{sr:eq:alpha-derivative}
\end{align}
These follow from the Hurwitz shift and derivative identities
\cite[25.11.3 and 25.11.17]{dlmf-hurwitz}. For any anchor with
$\Re\alpha_0>0$, a fixed-constant primitive is
\begin{align}
 \int_{\alpha_0}^{\alpha}D_0(t)\,dt
 ={}&(C-\mu)(\alpha-\alpha_0)
 +\lambda\log\frac{\Gamma(\alpha)}{\Gamma(\alpha_0)}\notag\\
 &+\zeta''(0,\alpha)-\zeta''(0,\alpha_0).
 \label{sr:eq:primitive}
\end{align}
The logarithm is the analytic log-Gamma on the right half-plane.
To verify the last term, differentiate
$\partial_\alpha\zeta(u,\alpha)=-u\zeta(1+u,\alpha)$ twice in
$u$ at zero; the result is
$\partial_\alpha\zeta''(0,\alpha)=2\gamma_1(\alpha)$.
This also explains why an anchored primitive involves the second
order derivative of Hurwitz zeta rather than an unspecified new constant.

\subsection{Two normalization checks and a rational specialization}
If every weight is multiplied by an arbitrary holomorphic $L(u)$
with $L(0)=1$, then $K_n$, $Q_n$, $\mathcal D$, and $\mathcal Z$
are all multiplied by $L(u)^2$. Therefore the cutoff at $u=0$ and
its constant $C$ are unchanged. Writing $l_1=L'(0)$ and
$l_2=L''(0)$, the spectral finite part becomes
\begin{equation}\label{sr:eq:normalization}
 \widetilde F=F+2l_1\lambda-l_1^2-l_2,
\end{equation}
whereas $\widetilde D_1=D_1+2l_1D_0$. These identities explicitly
track both amplitude derivatives relevant to a double pole.

At $a=c=1$ the weight is $w_n(0)=1/((n+1)(n+2))$ and
\[
 \mu=1,\quad \rho=\gamma-2,\quad \lambda=1-2\gamma,
 \quad C=\gamma-\gamma^2.
\]
An exact largest-index formula, obtained from the finite moments, is
\begin{equation}\label{sr:eq:rational-K}
 K_{v-1}(0)=1+\frac1{v+1}-\frac{H_v}{v}-\frac{H_v}{v+1},
 \qquad v\ge1.
\end{equation}
Its first value is zero, as required by the empty inner sum. Thus
\begin{equation}\label{sr:eq:integer-cutoff}
 \sum_{0\le m<n\le N}\frac{(n-m)^2}{(m+1)(m+2)(n+1)(n+2)}
 =N-\log^2N+(1-2\gamma)\log N+\gamma-\gamma^2+o(1).
\end{equation}
For the unnormalized weights the spectral finite part is
$F=\zeta(2)-2+2\gamma-2\gamma^2$. With the different normalization
$L(u)=\Gamma(1+u)$, the entire spectral product simplifies to
\[
 L(u)^2\mathcal Z(u)=\frac1{u^2(u-1)(1+u)^2},
\]
whose finite part is $-2$. Both have exactly the same finite-cutoff
constant $\gamma-\gamma^2$. This is a concrete second-order check
of \eqref{sr:eq:normalization}.

The specialization of \eqref{sr:eq:D0} at $a=c=\alpha=1$ yields
\[
 \sum_{v\ge1}\left\{
 \frac1{v+1}-\frac{H_v}{v}-\frac{H_v}{v+1}
 -\frac{1-2\gamma-2\log v}{v}\right\}
 =\gamma^2-1+2\gamma_1.
\]
This scalar specialization reduces to a classical harmonic--Stieltjes
relation; its novelty is not asserted. The general-shape theorem, the
uniform overlap subtraction, and the all-order cutoff comparison are
the substantive continuation of the proposed resonance problem.

\subsection{A functional polylogarithm identity and convergent endpoints}
The rational specialization has an exact disk identity as well as a
cutoff expansion. Define, for $|z|<1$,
\[
 \mathcal G(z)=\sum_{0\le m<n}
 \frac{(n-m)^2z^{n+1}}{(m+1)(m+2)(n+1)(n+2)}.
\]
The finite formula \eqref{sr:eq:rational-K} gives
\begin{equation}\label{sr:eq:G-polylog}
 \boxed{\displaystyle
 \mathcal G(z)=\frac z{1-z}-1-\frac{\log(1-z)}z
 -\operatorname{Li}_2(z)-\frac12\left(1+\frac1z\right)\log^2(1-z).}
\end{equation}
All apparent singularities at zero are removable. The first term is
$z^2/12$, corresponding to the configuration $(m,n)=(0,1)$.
To prove the formula, sum the four terms of
\eqref{sr:eq:rational-K}; the two harmonic transforms follow by
integrating $\sum_{v\ge1}H_vz^v=-\log(1-z)/(1-z)$ once, either with
$dz$ or $dz/z$. Thus this is a functional certificate before any
special-value substitution.

At $z=1/2$, the classical dilogarithm reflection identity
\cite[25.12.6]{dlmfpolylog} gives
\begin{equation}\label{sr:eq:half-polylog}
 \sum_{0\le m<n}\frac{(n-m)^2 2^{-(n+1)}}
 {(m+1)(m+2)(n+1)(n+2)}
 =2\log2-\log^22-\frac{\pi^2}{12}.
\end{equation}
For $p\ge2$, the ordinary convergent endpoint
\[
 J_p=\sum_{0\le m<n}\frac{(n-m)^2}
 {(m+1)(m+2)(n+1)(n+2)(n+1)^p}
\]
has a finite all-$p$ reduction. To distinguish multiple zeta values
from Hurwitz zeta, write
$\zeta_{\mathrm{MZ}}(j,1)=\sum_{v>m\ge1}v^{-j}m^{-1}$.
Set $E_j=\zeta_{\mathrm{MZ}}(j,1)+\zeta(j+1)$ for
$j\ge2$; Euler's height-one formula reduces every $E_j$ to products
of at most two ordinary zeta values \cite[Section~6]{incoming-determinantal}.
Then
\begin{equation}\label{sr:eq:allp-endpoint}
 J_p=\zeta(p)-E_{p+1}
 +\sum_{j=2}^{p}(-1)^{p-j}\bigl(\zeta(j)-E_j\bigr)
 +(-1)^{p-1}(1-\zeta(2)).
\end{equation}
Indeed partial fractions give
\begin{align*}
 \sum_{v\ge1}\frac1{v^p(v+1)}
 &=\sum_{j=2}^p(-1)^{p-j}\zeta(j)+(-1)^{p-1},\\
 \sum_{v\ge1}\frac{H_v}{v^p(v+1)}
 &=\sum_{j=2}^p(-1)^{p-j}E_j+(-1)^{p-1}\zeta(2);
\end{align*}
substitute them in \eqref{sr:eq:rational-K}. In particular,
\begin{align}
 J_2&=3\zeta(2)-2\zeta(3)-\tfrac54\zeta(4)-1,\label{sr:eq:J2}\\
 J_3&=1-2\zeta(2)+4\zeta(3)-\tfrac54\zeta(4)
       +\zeta(2)\zeta(3)-3\zeta(5).
 \label{sr:eq:J3}
\end{align}
The primitive of order one has a simple explicitly subtracted endpoint:
\begin{equation}\label{sr:eq:subtracted-J1}
 \sum_{v\ge1}\frac{K_{v-1}(0)-1}{v}
 =1-\zeta(2)-2\zeta(3).
\end{equation}
Its convergence follows from $K_{v-1}(0)-1=O(\log v/v)$.

\subsection{What this resolves, and what remains}
The next $r=2$ resonance requested in the earlier report is resolved:
two moment divergences require three asymptotic counterterm scales,
the overlap term is explicit, normal convergence is proved, and the
cutoff/spectral constants are compared including normalization changes.
The result does not establish a general higher-resonance theorem at
arbitrary particle number. Natural next problems are to identify the
escape strata at $r\ge3$, formulate an inclusion--exclusion subtraction
for simultaneous escape, and discover finite Pearson or contiguous
relations that reduce the number of independent divergent moments.
Another precise question is whether the order saving proved in
\eqref{sr:eq:lower-order-G} survives at higher particle number, where
several nontrivial residue amplitudes replace the single $R(u)$.

% Additive research section.  No source manuscript has been modified.
% Requires amsmath, amssymb, amsthm; theorem, proposition, corollary,
% remark environments.  Citation keys are left to the integrating article.
\section{All spectral coefficients from ordinary logarithmic Mellin moments}
\label{mellin:section}

The incoming \emph{Differential Polygamma and Lerch Jets} report
\cite[Section~9]{incoming-differential} asks for
spatial functionals selecting higher coefficients of its residue Dirichlet
germ, with the exact logarithmic normalization and endpoint constants.
The construction below answers that question at every order.  The
Mittag--Leffler expansion and the scalar residue filter are inherited
inputs; the shift-difference Mellin transform, its ordinary subtractions,
and the higher-coefficient formulas are the continuation proved here.
This is a claim of progress relative to the inspected material, not a
historical priority claim for all specializations of Mellin calculus.

\subsection{The data and the spatial normalization}

Let $P$ be a polynomial in finitely many variables $u_0,u_1,\ldots$ and
write
\[
 f_P(x)=P(\psi(x),\psi'(x),\ldots),\qquad
 \phi_P(Y)=P(Y,0,0,\ldots).
\]
Let its principal part at $-n$ be
\[
 \sum_{j=1}^{M}\frac{p_j(n)}{(x+n)^j},\qquad
 \Delta p_j(n)=p_j(n)-p_j(n-1),\qquad
 P_0(x)=\sum_{j=1}^{M}p_j(0)x^{-j}.
\]
Zero coefficients may be included.  These finite Laurent coefficients
are the harmonic polynomials supplied by the incoming report.  In
particular, $p_j(n)=O((1+\log n)^d)$ and
$\Delta p_j(n)=O((1+\log n)^d/n)$ for some fixed $d$.
Define
\begin{align}
 F_P(x)&=f_P(x)-f_P(x+1)-P_0(x),\label{mellin:F}\\
 c_P(L)&=\phi_P'(L)+p_1(0)=\sum_{r=0}^{R}c_rL^r,\label{mellin:c}\\
 h_P(x)&=F_P(x)+\mathbf 1_{[1,\infty)}(x)\frac{c_P(\log x)}x.
 \label{mellin:h}
\end{align}
The cutoff is exactly $1$, the same scale as the finite-part integral
on $[0,1]$.  A value assigned to the step function at $1$ has no effect.

\begin{theorem}[Ordinary Mellin realization of every residue coefficient]
\label{mellin:main}
The function $F_P$ extends analytically through $x=0$, and
\[
 h_P(x)=O((1+\log x)^d/x^2)\quad(x\longrightarrow\infty).
\]
Consequently
\begin{equation}\label{mellin:A}
 A_P(u)=\int_0^\infty x^{-u}h_P(x)\,dx
\end{equation}
is holomorphic in $-1<\operatorname{Re}u<1$.  Its derivatives there
are ordinary absolutely convergent integrals.  Put
\begin{equation}\label{mellin:B}
 \mathcal B_P(u)=\frac{\sin\pi u}{\pi}
 \left(A_P(u)-\sum_{r=0}^{R}\frac{c_r r!}{u^{r+1}}\right).
\end{equation}
This is a meromorphic germ at zero.  If
\[
 D_j(v)=\sum_{n\ge1}\frac{\Delta p_j(n)}{n^v},\qquad D_P=D_1,
\]
then the exact meromorphic identity is
\begin{equation}\label{mellin:bridge}
 \boxed{\mathcal B_P(u)=D_P(u)+
 \sum_{j=2}^{M}\frac{(u)_{j-1}}{(j-1)!}D_j(u+j-1).}
\end{equation}
Every series in the sum over $j\ge2$ is normally convergent on
$\operatorname{Re}u>-1$.  In the original strip
$0<\operatorname{Re}u<1$ one may equivalently write
\begin{equation}\label{mellin:unsubtracted}
 \mathcal B_P(u)=\frac{\sin\pi u}{\pi}
                  \int_0^\infty x^{-u}F_P(x)\,dx.
\end{equation}
\end{theorem}

\begin{proof}
Subtract the normalized Mittag--Leffler formula for $f_P(x+1)$ from
that for $f_P(x)$.  The constants cancel.  Reindexing a finite cutoff
and passing to its limit gives
\begin{equation}\label{mellin:partialfractions}
 F_P(x)=\sum_{n\ge1}\sum_{j=1}^{M}
               \frac{\Delta p_j(n)}{(x+n)^j}.
\end{equation}
Indeed the uncancelled last term is a principal part evaluated at
$x+N+1$, and tends to zero as
$O((1+\log N)^d/N)$.  The new series is normally convergent on
compact sets away from the negative integers.  In particular it is
analytic at zero.  This also follows directly from
$f_P-P_0$ and $f_P(x+1)$ being analytic there.

On the positive real axis, the usual differentiated digamma asymptotics
give
\[
 f_P'(x)=\frac{\phi_P'(\log x)}x
           +O((1+\log x)^d/x^2).
\]
Every monomial containing a positive polygamma derivative contributes
only to the error.  Integrating this estimate from $x$ to $x+1$, and
using $P_0(x)=p_1(0)/x+O(x^{-2})$, proves
\[
 F_P(x)=-\frac{c_P(\log x)}x
              +O((1+\log x)^d/x^2).
\]
Thus the subtraction in \eqref{mellin:h} proves all assertions about
$A_P$, including differentiation under the integral sign.

For $0<\operatorname{Re}u<1$ absolute convergence permits insertion
of \eqref{mellin:partialfractions}.  The elementary beta integral is
\[
 \int_0^\infty\frac{x^{-u}}{(x+n)^j}\,dx
 =n^{1-u-j}\frac{\Gamma(1-u)\Gamma(j-1+u)}{\Gamma(j)}.
\]
Its multiplication by $\sin(\pi u)/\pi$ is
$n^{1-u-j}(u)_{j-1}/(j-1)!$, including $j=1$.
The coefficient bound above makes the resulting sum absolutely
convergent.  Also
\[
 \int_1^\infty x^{-u-1}(\log x)^r\,dx=\frac{r!}{u^{r+1}}.
\]
These two computations prove \eqref{mellin:bridge} and
\eqref{mellin:unsubtracted} in the strip.  The $j\ge2$ series
converge locally normally for $\operatorname{Re}u>-1$; the scalar
residue filter continues $D_P$ meromorphically there.  Uniqueness of
continuation proves the germ identity.  No integral at a divergent
endpoint has been differentiated without its stated subtraction.
\end{proof}

\subsection{An explicit coefficient formula at every order}

For $m,r\ge0$ define the finite polynomial
\begin{equation}\label{mellin:E}
 E_{m,r}(L)=[u^r]e^{-Lu}\prod_{k=1}^{m}(1+u/k).
\end{equation}
Here an empty product is one.  The first three are
\[
 E_{m,0}=1,\qquad E_{m,1}=H_m-L,\qquad
 E_{m,2}=\frac{(H_m-L)^2-H_m^{(2)}}2.
\]
Thus these are finite harmonic polynomials, or equivalently Bell
polynomials with cumulants $H_m-L$ and
$(-1)^{k-1}(k-1)!H_m^{(k)}$ for $k\ge2$.

Write $\mu_k(P)=\int_0^\infty h_P(x)(\log x)^k\,dx$.
For $q\ge0$ put
\begin{align}
 B_q(P)={}&
 \sum_{\substack{k\ge0\\2k+1\le q}}
 \frac{(-1)^{k+q-2k-1}\pi^{2k}}
 {(2k+1)!(q-2k-1)!}\,\mu_{q-2k-1}(P)\notag\\
 &-\sum_{\substack{0\le r\le R\\q+r\ {
m even}}}
 c_r r!\frac{(-1)^{(q+r)/2}\pi^{q+r}}{(q+r+1)!}.
 \label{mellin:Bq}
\end{align}
The first sum is empty for $q=0$.  It is often clearer computationally
to take the coefficient directly in \eqref{mellin:B}; formula
\eqref{mellin:Bq} specifies every sign and factorial.

\begin{corollary}[Higher coefficients and the complete Stieltjes term]
\label{mellin:coefficients}
For every integer $q\ge1$,
\begin{equation}\label{mellin:Dq}
 \boxed{[u^q]D_P(u)=B_q(P)-
 \sum_{j=2}^{M}\sum_{n\ge1}
 \frac{\Delta p_j(n)}{(j-1)n^{j-1}}
 E_{j-2,q-1}(\log n).}
\end{equation}
The double sum is finite in $j$ and absolutely convergent in $n$.
If the inherited scalar residue polynomial is
\[
 R_P(L)=-\frac{\phi_P(L+i\pi)-\phi_P(L-i\pi)}{2i\pi},
 \qquad R_P'(L)=\sum_{r=0}^{S}a_rL^r,
\]
put
$\delta_n=\Delta p_1(n)-R_P'(\log n)/n$.
For every $q\ge0$ its entire ordinary-series counterpart is
\begin{equation}\label{mellin:stieltjes}
 \boxed{[u^q]D_P(u)=\frac{(-1)^q}{q!}
 \left(\sum_{r=0}^{S}a_r\gamma_{q+r}
                  +\sum_{n\ge1}\delta_n(\log n)^q\right).}
\end{equation}
The second series converges absolutely.  In particular, the same
scalar projection controls \emph{all} Stieltjes corrections.  If
$\phi_P=0$, every displayed $\gamma_{q+r}$ term vanishes.
\end{corollary}

\begin{proof}
Expanding $A_P$ in \eqref{mellin:B} gives \eqref{mellin:Bq}.
For $j\ge2$,
\[
 \frac{(u)_{j-1}}{(j-1)!}n^{-u}
 =\frac{u}{j-1}e^{-u\log n}\prod_{k=1}^{j-2}(1+u/k).
\]
Taking the coefficient of $u^q$ in the normally convergent identity
\eqref{mellin:bridge} proves \eqref{mellin:Dq}.  The logarithmic
factor grows by only a fixed power, so the convergence assertion
persists at every $q$.

The scalar filter gives
$\delta_n=O((1+\log n)^d/n^2)$ and
\[
 D_P(u)=\sum_r a_r(-1)^r\zeta^{(r)}(1+u)
                         +\sum_{n\ge1}\delta_nn^{-u}.
\]
With the convention
$\zeta(1+u)=u^{-1}+\sum_{m\ge0}(-1)^m\gamma_m u^m/m!$,
the coefficient of $u^q$ in
$(-1)^r\zeta^{(r)}(1+u)$ is
$(-1)^q\gamma_{q+r}/q!$.  This proves \eqref{mellin:stieltjes}.
\end{proof}

For $q=1$, the last sum in \eqref{mellin:Dq} is exactly the inherited
higher-pole correction
\[
 J_P=\sum_{j=2}^{M}\sum_{n\ge1}
              \frac{\Delta p_j(n)}{(j-1)n^{j-1}}.
\]
Hence $[u]\mathcal B_P=J_P+[u]D_P$, and the prior finite-part theorem
is recovered with its original endpoint constant:
\[
 \operatorname{FP}\int_0^1f_P(x)\,dx
       =C_P+[u]\mathcal B_P(u),\qquad
 C_P=\frac1{2\pi}\int_{-\pi}^{\pi}\phi_P(i\theta)\,d\theta.
\]
This is a normalization check, not a new proof of that antecedent.

\begin{corollary}[The first new logarithmic spatial functional]
\label{mellin:mixedq2}
If $P$ is mixed, meaning $\phi_P=0$, then
\begin{equation}\label{mellin:q2}
 \boxed{[u^2]D_P(u)=
 -\int_0^\infty h_P(x)\log x\,dx+\frac{\pi^2p_1(0)}6
 -\sum_{j=2}^{M}\sum_{n\ge1}
 \frac{\Delta p_j(n)(H_{j-2}-\log n)}{(j-1)n^{j-1}}.}
\end{equation}
All integrals and series are ordinary and absolutely convergent.
\end{corollary}

The local coefficient $p_1(0)$ is essential even in the mixed case.
It is the correction from $\sin(\pi u)/(\pi u)$, and its coefficient
is invisible at first order.  Thus disappearance of the Stieltjes
terms does not permit omission of this second-order endpoint term.

\subsection{Concrete polygamma and nonlinear examples}

For $P=u_1^2$ one has
\[
 p_4(n)=1,\quad p_2(n)=2(H_n^{(2)}+\zeta(2)),\quad
 p_1(n)=4(H_n^{(3)}-\zeta(3)).
\]
All omitted principal-part coefficients vanish.  Therefore
\begin{equation}\label{mellin:trigammasquare}
 \mathcal B_{u_1^2}(u)=(4+2u)\zeta(u+3).
\end{equation}
Let $Q(x)=\psi'(1+x)$.  An ordinary all-order generating formula is
obtained from \eqref{mellin:B} using
\[
 h_{u_1^2}(x)=
 \begin{cases}
 2\{Q(x)-\zeta(2)+2\zeta(3)x\}/x^2,&0<x<1,\\
 2\{Q(x)-\zeta(2)\}/x^2,&x>1.
 \end{cases}
\]
Its quadratic coefficient yields the explicit identity
\begin{equation}\label{mellin:trigamma-log}
 \boxed{\begin{aligned}
 &\int_0^1\frac{\psi'(1+x)-\zeta(2)+2\zeta(3)x}{x^2}\log x\,dx\\
 &\quad+\int_1^\infty\frac{\psi'(1+x)-\zeta(2)}{x^2}\log x\,dx\\
 &\qquad=-\zeta''(3)-\zeta'(3)-\frac{\pi^2\zeta(3)}3.
 \end{aligned}}
\end{equation}
Primes on $\zeta$ denote derivatives in the spectral order.  The
linear polygamma example provides a particularly simple independent
check of the endpoint term in \eqref{mellin:q2}.

For an example retaining a nonlinear spatial integrand, take $P=u_1^3$.
Write $Z(s,t)=\sum_{n>m\ge1}n^{-s}m^{-t}$ in its convergent region.
Finite Laurent extraction gives
\begin{align*}
 p_6&=1,&p_5&=0,&p_4&=3c_2,&p_3&=6c_3,\\
 p_2&=9c_4+3c_2^2,&p_1&=12c_5+12c_2c_3,
 \qquad c_r=H_n^{(r)}+(-1)^r\zeta(r).
\end{align*}
Inserting their exact differences into \eqref{mellin:bridge} proves
the finite depth-two identity
\begin{equation}\label{mellin:trigammacube}
 \boxed{\begin{aligned}
 \mathcal B_{u_1^3}(u)={}&
 \frac{u^3+9u^2+32u+48}{2}\,\zeta(u+5)\\
 &+6(u+2)\{Z(u+3,2)+\zeta(2)\zeta(u+3)\}\\
 &+12\{Z(u+2,3)-\zeta(3)\zeta(u+2)\}.
 \end{aligned}}
\end{equation}
All $Z$ functions in a neighborhood of zero are ordinary absolutely
convergent strict series.  Arbitrary higher order derivatives may be
taken there; no identity restricted to integer arguments is being
differentiated.

For clarity, put $K=\zeta(5)+\zeta(2)\zeta(3)$ and
\begin{align*}
 U(x)={}&\frac{3Q(x)}{x^4}+\frac{3Q(x)^2}{x^2}
 -\frac{3\zeta(2)}{x^4}+\frac{6\zeta(3)}{x^3}
 -\frac{9\zeta(4)+3\zeta(2)^2}{x^2}.
\end{align*}
The ordinary function for \eqref{mellin:A} is
\[
 h_{u_1^3}(x)=U(x)+\frac{12K}{x}\mathbf1_{(0,1)}(x).
\]
The displayed cancellations make it bounded at zero; for numerical
work its convergent Taylor series avoids loss of digits there.
With $\partial_1$ denoting a genuine first-order spectral derivative,
the second coefficient of \eqref{mellin:trigammacube} is
\begin{align}
 B_2(u_1^3)={}&12\zeta''(5)+16\zeta'(5)+\tfrac92\zeta(5)
 +6(\partial_1^2+\partial_1)Z(3,2)\notag\\
 &+6\zeta(2)\{\zeta''(3)+\zeta'(3)\}
 +6\partial_1^2 Z(2,3)-6\zeta(3)\zeta''(2).
 \label{mellin:cube-B2}
\end{align}
It follows that the following nonlinear logarithmic integral has a
finite, convergent depth-two evaluation:
\begin{equation}\label{mellin:cube-log}
 \boxed{\int_0^\infty h_{u_1^3}(x)\log x\,dx
                  =-B_2(u_1^3)-2\pi^2K.}
\end{equation}
This evaluates the logarithmic spatial functional in a specified
spectral family.  It makes no claim that the retained derivatives
of $Z$ are arithmetically irreducible.

\subsection{What this settles and what remains}

The higher-spectral-coefficient question is resolved for every
differential polynomial: the spatial kernel, the scale, all endpoint
terms, the higher-pole corrections, and every Stieltjes coefficient are
explicit.  The inherited convergent-word compiler applies to the
coefficients in \eqref{mellin:Dq}: it only acquires powers of
$\log n$, hence higher derivatives in the existing spectral arguments.
Thus these formulas introduce no additional kind of multiple-zeta
coordinate.  They do not settle relations between the resulting
evaluated coordinates.  A separate next problem is to extend the
construction to independent shifts, where the pole lattices no longer
align under the unit difference, and to characterize which linear
combinations eliminate all higher-depth spectral derivatives.

\section{Audit conclusions, integration changes, and further research}
\label{sec:audit-research}

\subsection{The precise change in mathematical status}

The inspection covered the source statements used by this investigation,
their nearby conjectures, and the six incoming reports at the pinned
revision. It was not a line-by-line certification of the entire collective
manuscript. No new counterexample or invalid proof in those inherited
statements was established. Several tempting shortcuts in the present
derivations were rejected or corrected before delivery. The appropriate
repository changes are the following additions and status refinements.

\begin{longtable}{@{}>{\raggedright\arraybackslash}p{.23\textwidth}>{\raggedright\arraybackslash}p{.34\textwidth}>{\raggedright\arraybackslash}p{.36\textwidth}@{}}
\toprule
Target in the inspected material & Proved advancement & Boundary retained\\
\midrule
\endhead
All reconstruction rows, Bernoulli report R2 &
Theorem~\ref{thm:all-rows} gives every entry of $M_b^{-1}$;
Theorems~\ref{thm:extraction-series} and \ref{thm:polylog-filter}
give its ordinary and polylogarithmic evaluations. &
The inverse is for the $(2,2b)$ matrix. Formal coordinate rank is
not arithmetic independence.\\[.5em]
Optimal-prefix conjecture beyond $p=6$ &
Theorem~\ref{rank8:theorem} proves the $p=8$, $q=2b$ family
for $b\ge5$, $b\not\equiv1\pmod {16}$. &
The missing residue class and general unequal orders $p\ge10$
are not settled by this paper.\\[.5em]
Centered cubic tails &
Theorem~\ref{cc:all} evaluates every regular odd centered cubic
trigamma moment; Theorems~\ref{cc:generator} and \ref{cc:polylog}
give its analytic and finite polylogarithmic generators. &
The lowest cubic evaluation has classical antecedents.
No general higher-power closure theorem is claimed.\\[.5em]
Next two-particle resonance &
Theorems~\ref{sr:thm:germ} and \ref{sr:thm:constant} establish
the grouped subtraction and exact cutoff constant. Corollary
\ref{sr:cor:alljets} gives every spectral row. &
The result concerns two particles and the specified next resonance.
Arbitrary particle number and arbitrary resonance remain open.\\[.5em]
Higher spatial spectral functionals &
Theorem~\ref{mellin:main} and Corollary~\ref{mellin:coefficients}
give ordinary logarithmic Mellin formulas at every coefficient order. &
The statement uses the inherited unshifted differential-polynomial
residue germ. General independent spatial shifts need a new analysis.\\[.5em]
Short Gaussian $S_6$ and revised $S_8$ &
No status change. &
The current frozen candidates remain conjectural. The older rejected
$S_8$ vector is a different statement.\\
\bottomrule
\end{longtable}

For integration, the full inverse belongs immediately after the rational
row in the Bernoulli moment-rank development. The rank-eight theorem is
an additional case of the existing optimal-prefix programme, and should
be attached with its congruence qualification. The cubic section belongs
after the canonical centered-square chapter. The resonance theorem
continues the determinantal report; the Mellin theorem continues the
differential-polygamma report. These are additive changes. The package's
integration notes identify the source files and supplied artifacts.

\subsection{Normalization issues that must survive editing}

\paragraph{The omitted binomial term.}
The lower-index-zero summand in the source tail matrix is not present.
Restoring it changes the matrix, the exceptional zeroth column, and the
reconstruction coefficients. Formula~\eqref{eq:M-system} and the direct
finite-cutoff verification in the package fix this convention explicitly.

\paragraph{The zeroth power in Faulhaber's formula.}
The uniformly correct expression for $\sum_{k=1}^N k^m$, $m\ge0$,
is \eqref{eq:faulhaber-tail}, with $B_{m+1}(1)$ at the lower endpoint.
Replacing it by $B_{m+1}(0)$ is valid for $m\ge1$ but fails at $m=0$.
This scope error was found in an early draft of the present article and
corrected. It is not reported as an error in an inherited theorem.

\paragraph{Exceptional cubic boundary terms.}
The $-\zeta(2)/2$ term at $j=2$ and the rational $1/4$ term at $j=3$
in \eqref{cc:Cj} are required. The latter is the surviving
$B_1^2/B_1^3$ boundary contribution in \eqref{cc:boundaryexception}.
Deleting low-order exceptions because a generic parity pattern looks
simpler would change the convergent identities.

\paragraph{Cutoff constants and spectral finite parts.}
At the second resonance, $C$ in \eqref{sr:eq:lambda-C} and $F$ in
\eqref{sr:eq:spectral-laurent} are different precisely normalized
quantities. Equation~\eqref{sr:eq:normalization} includes both derivatives
of the amplitude needed at a double pole. The example $a=c=1$ gives
the same cutoff constant $\gamma-\gamma^2$ but different spectral
finite parts under two normalizations. This is a proved comparison,
not an assertion that the inspected report made the incorrect identification.

\paragraph{A second-order endpoint constant after scalar cancellation.}
When $\phi_P=0$, the scalar-filter Stieltjes terms disappear, but
$\pi^2p_1(0)/6$ remains in \eqref{mellin:q2}. It comes from the sine
factor in the Mellin transform. Its absence at first order provides
no reason to discard it at second order.

\paragraph{An extrapolation that is actually false.}
The rank-eight determinant's valuation follows a simple law on the
proved residue classes, but the unrestricted extrapolation fails at
$b=65$ and $129$, as documented in the rank-eight section. These are
counterexamples to a guess considered during this investigation, not
to the source optimal-prefix conjecture. The determinants at those
points are nonzero. The missing class must be addressed by a different
argument or a finer congruence analysis.

\subsection{Exact certificates and independent diagnostics}

The accompanying scripts separate identities over $\Q$ from numerical
comparisons. The reconstruction verifier checks both products with the
original binomial matrix through $b=15$, the shifted and unshifted
orthogonality identities, the exceptional inverse row, and the printed
examples. A second implementation reconstructs 35 late-coordinate rows
directly from finite summation and the uncentered Hurwitz expansion;
it does not insert the proposed matrix formula. Exact coefficient tables
contain all 20 extraction rows for $2\le b\le6$.

The rank-eight computation is a finite symbolic certificate in an
integer polynomial ring, followed by a proved congruence for the Bernoulli
inputs. Checking all coefficients of that polynomial identity is what
makes it apply to an infinite family. A list of specialized nonzero
determinants would not suffice. The full polynomial and a generator that
reconstructs it from the four-by-four matrix are supplied.

The cubic verifier independently reconstructs 18 coordinates from finite
cutoffs, checks the generating-function coefficients, and differentiates
the finite polylogarithmic numerator twice with exact zero residual.
It also checks the rational and $\zeta(3)$ extraction identities.
The resonance verifier checks the finite Pearson relations and the
Laurent cancellations algebraically, and evaluates defining finite
weight sums with an independently generated asymptotic tail. The Mellin
verifier checks 56 exact principal-part and coefficient identities,
then compares spatial quadrature with spectral expressions.

Numerical diagnostics use 50, 80, or 110 decimal working digits, as
recorded by each script. The six reconstruction comparisons have largest
absolute residual below $3.4\times10^{-59}$; the cubic moment checks
through $r=8$ have largest higher-cutoff residual below
$4.8\times10^{-76}$. The two nonzero-parameter resonance comparisons,
including a complex parameter, agree within $1.4\times10^{-30}$.
The nine Mellin comparisons have largest relative residual below
$1.1\times10^{-49}$. These are independently organized diagnostics,
not rigorous interval enclosures. Cancellation can exhaust working
precision at large polynomial degree, so the reports retain the cutoff,
expansion order, and successive comparison rather than a single
unqualified count of ``correct digits.'' The analytic proofs establish
the infinite-family statements.

\subsection{Further research questions}

The following questions retain a concrete exact target. They prioritize
identities and reductions; estimates are tools for justifying the
manipulations, rather than separate objectives.

\begin{enumerate}[leftmargin=*,label=\textbf{Q\arabic*.}]

\item \textbf{Finish rank eight on $b\equiv1\pmod {16}$.}
Determine an arithmetic invariant of the four-by-four perturbation that
cannot vanish in this residue class. Higher modulus versions of the
Faulhaber congruence are available, but the counterexamples to the simple
valuation law show that its exponent cannot be predicted merely from
$v_2(b-1)$. A useful result would prove nonvanishing for all remaining
$b$, or identify a complete, exact valuation recursion with a termination
argument. Checking more values alone does not settle the question.

\item \textbf{Prove the optimal prefix for general unequal even orders.}
For $(p,q)=(2a,2b)$ with $a\ge5$, the inherited reduction leaves an
$a$-square perturbation determinant. Seek a structural factorization,
a congruence stable in $a$, or an algebraic interpretation of its
nonvanishing. The Bernoulli functional is indefinite, so a positive
moment-matrix argument cannot simply be imported. The result should
identify the exceptional arithmetic factors, not leave their nonvanishing
as an implicit assumption.

\item \textbf{Construct all higher-$p$ reconstruction rows.}
The kernel of Theorem~\ref{thm:kernel} and a finite-rank inverse formula
suggest a route for $(4,2b)$ and $(6,2b)$, whose rank is already known.
The target is an explicit coefficient formula comparable to
\eqref{eq:all-rows}, with no unknown minors. Translate those coefficients
into ordinary polynomial-subtracted polygamma sums with one prescribed
even-zeta coordinate. The independent finite-cutoff method in the package
can verify their normalization before an all-order proof is attempted.

\item \textbf{Extend the reconstruction to shifted and character tails.}
Replace $A_p(n)$ by a Hurwitz tail or a Dirichlet-character tail.
Determine the correct generalized Bernoulli kernel and state the exact
finite correction when the midpoint identity
$\zeta(-2j,1/2)=0$ is lost. One can first ask for a single rational
shift and a common character, then for distinct shifts. A claimed
ordinary evaluation must include the changed subtraction constants.

\item \textbf{Classify higher powers of the centered trigamma tail.}
For $d\ge4$, examine
$\sum_{n\ge0}\psi'(n+1)^d x_n^{-s}$ at its regular integer values.
Does the correct parity family lie in a fixed finite coordinate space
for each $d$? The cubic proof identifies a finite cancellation of linear
tails before taking constant terms. Find its higher-power analogue, or
prove a specific obstruction. An explicit quartic formula and a generator
would be a decisive next case; no quartic closure is assumed here.

\item \textbf{Determine the functional depth of the cubic generator.}
Theorem~\ref{cc:polylog} uses six specified depth-two polylogarithms.
Determine whether their particular combination admits a smaller function
basis under an explicitly stated class of argument transformations.
An equality of derivatives together with an anchored value is a suitable
certificate. A numerical relation at one $e^{-t}$ does not establish a
functional reduction, and formal symbol rank alone does not establish
arithmetic irreducibility at an endpoint.

\item \textbf{Integrate the new boundary operators at interior arguments.}
The operators in \eqref{eq:Phi-filter} prescribe exact cancellations at
$z=1$. Find anchored primitives that lower their differential order and
give compact formulas at interior algebraic points, including $1/2$ and
cyclotomic points accessible by continuation. Track all integration
constants from the power series at zero. The rational resonance generator
in \eqref{sr:eq:G-polylog} supplies a simple model of this approach.

\item \textbf{Develop a general simultaneous-escape subtraction.}
For three or more particles, classify the regions in which subsets of
summation indices tend to infinity together. Determine an exact
inclusion--exclusion subtraction and the associated exponents in the
Hurwitz terms. The doubled scale $1+2u$ in the present two-particle
case is the first nontrivial example. A successful theorem must prove
normal convergence of the full grouped expression before differentiating
at a resonance.

\item \textbf{Explain the saving of one shape-polygamma order.}
Equation~\eqref{sr:eq:lower-order-G} removes the apparent $R_{k+2}$
dependence from every $k$th spectral row. Is this a consequence of a
general determinant identity or of the two-particle Pearson reduction?
Find a corresponding regular expression at higher resonances or prove
where the saving fails. Such a result would reduce the number of
independent harmonic weights needed in exact derivative formulas.

\item \textbf{Extend the ordinary Mellin bridge to independent shifts.}
For $P$ in polygamma functions evaluated at several distinct shifts,
formulate the exact shift-difference principal parts and identify the
replacement for $D_j(u+j-1)$. A natural target is a normally convergent
Lerch or Hurwitz multiple-zeta family with all endpoint logarithms
explicit. The first problem is to specify the spatial subtraction scale
and its transformation law; the generalized Stieltjes corrections then
have an unambiguous meaning.

\item \textbf{Reduce the concrete nonlinear spectral derivatives.}
The cube identity \eqref{mellin:trigammacube} provides particular
combinations of $\partial_1^j Z(s,t)$ at convergent integer pairs.
Determine which combinations reduce to ordinary zeta derivatives,
height-one coefficients, or the existing Tornheim coordinates. Functional
relations must hold in a neighborhood of the differentiated point;
an identity true only at integer orders cannot be differentiated as if
it were such a relation. The logarithmic integral
\eqref{mellin:cube-log} is a concrete initial target.

\item \textbf{Compute the spectral image of the differential quotient.}
The incoming report classifies integration by parts in a free
differential-polynomial algebra. The Mellin bridge supplies a map from
that quotient, with its endpoint functional, to harmonic residue series.
Determine the kernel of this map at fixed weight, first exactly at low
weight and then by structural generators. A formal quotient dimension
and a dimension after substituting the digamma function need not agree.

\item \textbf{Build anchored multivariable primitive systems.}
Combine the spectral derivatives in \eqref{sr:eq:alljets} with the
shift primitive \eqref{sr:eq:primitive}. Determine finite recurrences for
iterated primitives in the shift and shape parameters, with base values
fixed at specified anchors. The goal is an identity in known Gamma,
Hurwitz-zeta, Stieltjes, and multiple-zeta families, rather than naming
an undetermined constant after each integration.

\item \textbf{Return to the frozen Gaussian $S_6$ and revised $S_8$ targets.}
Preserve their exact current vectors and the repository's established
shuffle equivalences. Seek a functional or iterated-integral relation
whose rigorously evaluated boundary is precisely one of those vectors.
The present identities expand the available analytic tools but establish
no direct implication for those two conjectures. Their status should
change only after such a relation and its endpoint evaluation are proved.

\end{enumerate}

The concrete deliverable of this investigation is therefore a set of
proved reconstruction, moment, resonance, and Mellin theorems with
explicit coefficients and normalization constants. The unresolved
questions are kept attached to the exact missing mathematical step.

\begin{thebibliography}{99}
\raggedright
\addcontentsline{toc}{section}{References and inspected sources}

\bibitem{repo}
V. Reshetnikov and the ProveIt project,
\emph{Polylogarithms manuscript and incoming research reports},
repository revision \repoSHA, inspected 11 October 2026.
\href{https://github.com/VladimirReshetnikov/ProveIt/tree/446dca9bd65c20e409d1d2561d4ffad408103e2a/Analysis/Polylogarithms/docs/manuscript}{Pinned manuscript};
\href{https://github.com/VladimirReshetnikov/ProveIt/tree/446dca9bd65c20e409d1d2561d4ffad408103e2a/docs/incoming}{pinned incoming directory}.

\bibitem{incoming-bernoulli}
ProveIt incoming report,
\emph{Bernoulli, Tornheim, and digamma research}, 11 October 2026.
Archive \texttt{ProveIt\_Bernoulli\_Tornheim\_Digamma\_2026-10-11.zip};
especially the tail foundation, moment-rank, general-Hankel,
rank-four, rank-six, and audit/research sections.
\href{https://github.com/VladimirReshetnikov/ProveIt/blob/446dca9bd65c20e409d1d2561d4ffad408103e2a/docs/incoming/ProveIt_Bernoulli_Tornheim_Digamma_2026-10-11.zip}{Archive at inspected revision}.

\bibitem{incoming-cubic}
ProveIt incoming report,
\emph{Cubic tail descent}, 10 October 2026.
Archive \texttt{ProveIt\_Cubic\_Tail\_Descent\_2026-10-10.zip}.
\href{https://github.com/VladimirReshetnikov/ProveIt/blob/446dca9bd65c20e409d1d2561d4ffad408103e2a/docs/incoming/ProveIt_Cubic_Tail_Descent_2026-10-10.zip}{Archive at inspected revision}.

\bibitem{incoming-determinantal}
ProveIt incoming report,
\emph{Determinantal harmonic resonance}, 10 October 2026.
Archive \texttt{ProveIt\_Determinantal\_Harmonic\_Resonance\_2026-10-10.zip};
especially the Gauss determinant, the first resonance, the harmonic
endpoint reductions, and the next-resonance question in Section~9.
\href{https://github.com/VladimirReshetnikov/ProveIt/blob/446dca9bd65c20e409d1d2561d4ffad408103e2a/docs/incoming/ProveIt_Determinantal_Harmonic_Resonance_2026-10-10.zip}{Archive at inspected revision}.

\bibitem{incoming-differential}
ProveIt incoming report,
\emph{Differential polygamma and Lerch jets}, 11 October 2026.
Archive \texttt{ProveIt\_Differential\_Polygamma\_and\_Lerch\_Jets\_2026-10-11.zip};
especially the normalized Mittag--Leffler theorem, scalar residue
filter, and higher-spectral-derivative question in Section~9.
\href{https://github.com/VladimirReshetnikov/ProveIt/blob/446dca9bd65c20e409d1d2561d4ffad408103e2a/docs/incoming/ProveIt_Differential_Polygamma_and_Lerch_Jets_2026-10-11.zip}{Archive at inspected revision}.

\bibitem{incoming-multishift}
ProveIt incoming report,
\emph{Multishift reflection}, 10 October 2026.
Archive \texttt{ProveIt\_Multishift\_Reflection\_2026-10-10.zip}.
\href{https://github.com/VladimirReshetnikov/ProveIt/blob/446dca9bd65c20e409d1d2561d4ffad408103e2a/docs/incoming/ProveIt_Multishift_Reflection_2026-10-10.zip}{Archive at inspected revision}.

\bibitem{incoming-tail}
ProveIt incoming report,
\emph{Polygamma tail identities}, 11 October 2026.
Archive \texttt{ProveIt\_Polygamma\_Tail\_Identities\_2026-10-11.zip}.
\href{https://github.com/VladimirReshetnikov/ProveIt/blob/446dca9bd65c20e409d1d2561d4ffad408103e2a/docs/incoming/ProveIt_Polygamma_Tail_Identities_2026-10-11.zip}{Archive at inspected revision}.

\bibitem{cc:repo-square}
ProveIt manuscript,
\emph{Centered trigamma moments}, chapter source
\texttt{07-centered-trigamma-moments.tex}, at the revision of \cite{repo}.
\href{https://github.com/VladimirReshetnikov/ProveIt/blob/446dca9bd65c20e409d1d2561d4ffad408103e2a/Analysis/Polylogarithms/docs/manuscript/chapters/07-centered-trigamma-moments.tex}{Pinned chapter source}.

\bibitem{cc:furdui2015}
O. Furdui,
The evaluation of a quadratic and a cubic series with trigamma function,
\emph{Applied Mathematics E-Notes} \textbf{15} (2015), 187--196.
\href{https://www.math.nthu.edu.tw/~amen/2015/2015.htm}{Journal volume contents};
\href{https://www.kurims.kyoto-u.ac.jp/EMIS/journals/AMEN/2015/AMEN(150526).pdf}{journal article mirror}.

\bibitem{cc:hoffman2016}
M. E. Hoffman,
Sums of products of Riemann zeta tails,
\emph{Mediterranean Journal of Mathematics} \textbf{13} (2016),
2771--2781.
\href{https://doi.org/10.1007/s00009-015-0653-9}{doi:10.1007/s00009-015-0653-9};
\href{https://arxiv.org/abs/1602.03399}{arXiv:1602.03399}.

\bibitem{dilcher-jiu}
K. Dilcher and L. Jiu,
\emph{Hankel determinants of sequences related to Bernoulli and Euler
polynomials}, arXiv:2007.09821 (2020).
\url{https://arxiv.org/abs/2007.09821}.
Used as background for the Hankel and orthogonal-polynomial method;
the particular kernel formulas in this article are proved explicitly.

\bibitem{dlmf-op}
NIST Digital Library of Mathematical Functions,
\S18.2, General orthogonal polynomials.
\url{https://dlmf.nist.gov/18.2}.

\bibitem{dlmf-hurwitz}
NIST Digital Library of Mathematical Functions,
\S25.11, Hurwitz zeta function.
\url{https://dlmf.nist.gov/25.11}.

\bibitem{dlmf-polygamma}
NIST Digital Library of Mathematical Functions,
\S5.15, Polygamma functions.
\url{https://dlmf.nist.gov/5.15}.

\bibitem{dlmfgauss}
NIST Digital Library of Mathematical Functions,
\S15.4, Hypergeometric special cases, especially (15.4.20).
\url{https://dlmf.nist.gov/15.4}.

\bibitem{dlmfgammaratio}
NIST Digital Library of Mathematical Functions,
\S5.11, Gamma asymptotic expansions, especially (5.11.13).
\url{https://dlmf.nist.gov/5.11}.

\bibitem{dlmfpolylog}
NIST Digital Library of Mathematical Functions,
\S25.12, Polylogarithms, especially (25.12.6).
\url{https://dlmf.nist.gov/25.12}.
\end{thebibliography}

\end{document}
